\chapter{前馈神经网络}
\label{chap:feedforward-networks}

第~\ref{chap:neural-network-overview}章用异或说明了隐藏表示的作用：两个非线性单元就能把原始坐标中无法线性分开的样本变成可分的表示。但那组权重是手工给出的。面对真实数据，网络应当怎样组织这些单元，训练目标怎样决定各层参数的变化，又凭什么相信训练结果可以用于未见样本？这三个问题分别关系到计算结构、参数学习和预测能力，不能仅凭网络层数或训练误差回答。

本章以层状、全连接的前馈网络为主要对象，把前章的函数复合落实为可计算、可求导的模型。网络结构规定候选函数及其输出含义，参数学习把最终损失对各层参数的影响算出来，理论分析则区分表达一个函数与从有限数据学得低风险预测规则。常见应用把这些要素接回完整的输入、训练和预测过程。读者可以沿用前文的向量、矩阵、损失和梯度知识；复杂度与稳定性的结论主要承接
\BookChapterRef{chap:deep-learning-theory}{第一卷“基础、理论和可信性”的“深度学习的可学习性”章}，
不在这里重复一般证明。

\section{网络结构}
\label{sec:ffn-structure}

前馈性质规定的是计算依赖的方向：同一次前向计算中，后续结果不会沿有向环返回上游。它并没有要求每个单元必须连接所有输入，也没有要求所有层宽相同。本节采用最直接的全连接结构，说明一次预测究竟计算什么，再考察激活函数和输出目标怎样改变这种计算。

\subsection{神经元}
\label{sec:ffn-neuron}

要把一组输入变成新的特征，可以让一个人工神经元读取若干数值，计算加权和，再通过激活函数得到输出。设单元的输入为$\bm{u}\in\mathbb{R}^{m}$，权重为$\bm{w}\in\mathbb{R}^{m}$，偏置为$b\in\mathbb{R}$，则
\begin{equation}
  z=\bm{w}^{\top}\bm{u}+b,
  \qquad h=\sigma(z)\text{。}
  \label{eq:ffn-unit}
\end{equation}
权重决定各个输入怎样共同影响单元，偏置改变单元响应的位置。这里的$z$称为\term{预激活值}（pre-activation），即非线性变换之前的加权结果；$h$是单元输出。没有偏置时，ReLU单元的分界$\bm{w}^{\top}\bm{u}=0$必须经过原点，加入偏置后便能移动这条分界。

同一层的多个单元可以各自读取整组输入，用不同权重计算不同坐标。这种\term{全连接层}（fully connected layer）让下一层每个单元与上一层每个单元相连。把一个单元的权重按行放入矩阵，就能用一次矩阵运算计算整层。沿用前章的列向量约定，设网络有$L$个隐藏层，输入为$\bm{x}\in\mathbb{R}^{d}$，令$d_0=d$、$\bm{h}^{(0)}=\bm{x}$，各隐藏层依次计算
\begin{equation}
  \begin{aligned}
    \bm{z}^{(\ell)}
      &=\bm{W}_\ell\bm{h}^{(\ell-1)}+\bm{b}_\ell,\\
    \bm{h}^{(\ell)}&=\sigma_\ell\!\left(\bm{z}^{(\ell)}\right),
    \qquad \ell=1,\ldots,L\text{。}
  \end{aligned}
  \label{eq:ffn-hidden-layers}
\end{equation}
其中$\bm{W}_\ell\in\mathbb{R}^{d_\ell\times d_{\ell-1}}$、$\bm{b}_\ell\in\mathbb{R}^{d_\ell}$，$\bm{z}^{(\ell)}$和$\bm{h}^{(\ell)}$都属于$\mathbb{R}^{d_\ell}$，$\sigma_\ell$逐坐标作用。上标$\ell$表示层的位置，不是幂，也不是参数更新的迭代步。

隐藏层之后还需一个与任务相配的输出层。记原始输出为$\bm{a}\in\mathbb{R}^{q}$、$d_{L+1}=q$，完整模型写为
\begin{equation}
  \begin{aligned}
    \bm{a}&=\bm{W}_{L+1}\bm{h}^{(L)}+\bm{b}_{L+1},\\
    \widehat{\bm{y}}&=g(\bm{a})\text{。}
  \end{aligned}
  \label{eq:ffn-output-map}
\end{equation}
矩阵$\bm{W}_{L+1}\in\mathbb{R}^{q\times d_L}$和向量$\bm{b}_{L+1}\in\mathbb{R}^{q}$参与学习，输出变换$g$由任务决定。$g$可以是恒等映射，也可以把得分转换为概率。用$\bm{\theta}\in\mathbb{R}^{p}$收集全部权重与偏置，模型记为$\bm{f}_{\bm{\theta}}(\bm{x})=\widehat{\bm{y}}$；标量输出时写$f_{\bm{\theta}}$。

图~\ref{fig:ffn-architecture}展示了这种分层依赖。输入层只提供数据，不计算一组新的可学习参数；隐藏层的每个坐标都是一次式\eqref{eq:ffn-unit}的计算。通常称这种具有一个或多个全连接隐藏层的网络为\term{多层感知机}（multilayer perceptron, MLP），即用多层非线性变换构造表示的前馈网络。这个名称不意味着现代MLP必须使用早期感知机的硬阈值激活。

\begin{figure*}[htb]
  \centering
  % 层宽 3,4,3,2。浅色神经元、细连接与少量方向提示构成统一视觉语法。
\begin{tikzpicture}[x=1mm,y=1mm,font=\figurefont\footnotesize]
  \foreach \i/\y in {1/28,2/14,3/0}
    \coordinate (x\i) at (12,\y);
  \foreach \i/\y in {1/35,2/21,3/7,4/-7}
    \coordinate (h\i) at (52,\y);
  \foreach \i/\y in {1/28,2/14,3/0}
    \coordinate (k\i) at (92,\y);
  \foreach \i/\y in {1/21,2/7}
    \coordinate (a\i) at (132,\y);
  \foreach \i in {1,2,3} \foreach \j in {1,2,3,4}
    \draw[nn link,draw=NNInput!58,shorten <=4.4mm,shorten >=4.4mm] (x\i)--(h\j);
  \foreach \i in {1,2,3,4} \foreach \j in {1,2,3}
    \draw[nn link,draw=NNHidden!58,shorten <=4.4mm,shorten >=4.4mm] (h\i)--(k\j);
  \foreach \i in {1,2,3} \foreach \j in {1,2}
    \draw[nn link,draw=NNOutput!58,shorten <=4.4mm,shorten >=4.4mm] (k\i)--(a\j);
  \foreach \i in {1,2,3}
    \node[nn input] at (x\i) {$x_{\i}$};
  \foreach \i in {1,2,3,4}
    \node[nn hidden] at (h\i) {$h_{\i}$};
  \foreach \i in {1,2,3}
    \node[nn hidden] at (k\i) {$h_{\i}$};
  \foreach \i in {1,2}
    \node[nn output] at (a\i) {$a_{\i}$};
  \foreach \x/\title/\symbol in {
    12/输入层/\bm{x},52/隐藏层1/\bm{h}^{(1)},
    92/隐藏层2/\bm{h}^{(2)},132/输出层/\bm{a}}{
    \node[nn heading] at (\x,55) {\title};
    \node[text=BookMuted] at (\x,47) {$\symbol$};
  }
  \draw[nn flow] (26,47)--(38,47);
  \draw[nn flow,draw=NNHidden] (66,47)--(78,47);
  \draw[nn flow,draw=NNOutput] (106,47)--(118,47);
  \node[nn annotation] at (72,-19)
    {$\bm{z}^{(\ell)}=\bm{W}_{\ell}\bm{h}^{(\ell-1)}+\bm{b}_{\ell}$
     \qquad $\bm{h}^{(\ell)}=\sigma_{\ell}(\bm{z}^{(\ell)})$};
  \node[nn annotation] at (132,-7) {$\widehat{\bm{y}}=g(\bm{a})$};
  % 无底色、无外框图例，沿用函数图的短样线与文本排列。
  \foreach \x/\col/\txt in {5/NNInput/输入坐标,57/NNHidden/隐藏响应,109/NNOutput/原始得分}{
    \draw[\col,line width=.9pt] (\x,-32)--++(10,0);
    \node[circle,draw=\col,fill=\col!16,line width=1.1pt,
      minimum size=3.6mm,inner sep=0pt] at (\x+5,-32) {};
    \node[anchor=west] at (\x+13,-32) {\txt};
  }
\end{tikzpicture}

  \caption{两个隐藏层的全连接前馈网络。蓝色圆点提供输入$\bm{x}$，紫色圆点计算隐藏表示，青色圆点产生原始得分$\bm{a}$，之后再使用任务规定的输出变换$g$。细线表示可学习的连接，层间箭头标明前向计算方向；每个目标单元还具有一个未画出的偏置。图中层宽依次为$3,4,3,2$，只是结构示例。}
  \label{fig:ffn-architecture}
\end{figure*}

对没有参数共享的上述网络，参数总数为
\begin{equation}
  p=\sum_{\ell=1}^{L+1}d_\ell(d_{\ell-1}+1)\text{。}
  \label{eq:ffn-parameter-count}
\end{equation}
括号内的两项分别来自每个单元的输入权重和偏置。若输入维度为$100$，两个隐藏层宽度均为$64$，输出维度为$1$，参数量为$64\times101+64\times65+65=10689$。这说明输入编码和层宽会直接改变模型规模。忽略激活和偏置的较小开销，一次单样本前向计算的乘加量与$\sum_\ell d_\ell d_{\ell-1}$同阶；参数量并不能单独说明表达是否高效或预测是否可靠。

本文把$L$称为隐藏层数，把$d_\ell$称为第$\ell$层宽度。文献中的\term{深度}（depth）有时包含输出层，有时只计隐藏层，比较结果时应核对口径。跨层连接只要不构成有向环，也可以保持前馈性质；本章的逐层推导采用普通串联结构。

\subsection{激活函数}
\label{sec:ffn-activations}

前章已经证明，连续叠加仿射层仍然只是一次仿射变换。隐藏层需要非线性激活，使不同输入范围可以有不同响应。但非线性也会改变导数：有些函数把大幅输入变化压成很小的输出变化，有些函数把一部分输入直接置零，还有些函数让响应形状参与学习。选择激活函数时，需要同时考察表示、梯度、输出约束和计算成本。

本节按响应方式介绍常见函数：恒等与阶跃作为基准，双侧饱和函数限制输出范围，整流与指数型函数保留较大的正响应，自门控函数平滑地调节输入，Maxout与Softmax则联合处理多个预激活值。带参数函数的图像必须固定参数才能比较；随机函数还需区分训练与预测阶段。除特别说明外，输入记为标量$z$，输出记为$\sigma(z)$。

\subsubsection{恒等与阶跃}
\label{sec:ffn-activation-baselines}

若需要保留任意实数输出，可以使用\term{恒等函数}（identity function）；若只需要一个硬的二值判定，可以使用\term{阶跃函数}（step function）。两者分别为
\begin{equation}
  \operatorname{Identity}(z)=z,
  \qquad
  H(z)=\begin{cases}0,&z<0,\\1,&z\geq0\text{。}\end{cases}
  \label{eq:ffn-identity-step}
\end{equation}
恒等函数的导数为$1$，常用于回归输出或无需额外激活的投影；全部隐藏层都用恒等函数时，多层串联仍是仿射映射。阶跃函数适合表达阈值决策，但除阈值处不可微之外，在其余位置导数都为$0$，普通反向传播无法据此学习阈值单元的输入参数。预测时作硬判定与训练时使用连续可导替代，是两种不同用途。

\begin{figure*}[htb]
  \centering
  % 原创解析曲线。采样数据由 feedforward-curves.py 生成。
\begingroup
% 局部统一的科学绘图样式；曲线保持颜色和线型双重编码。
\pgfplotsset{
  ffn curve/.style={line width=1.3pt,line cap=round,line join=round},
  ffn reference curve/.style={NNInk!58,line width=0.95pt},
  ffn activation axis/.style={
    width=72mm,height=59mm,anchor=south west,
    axis lines=middle,
    axis line style={NNInk!65,line width=0.65pt},
    grid=major,grid style={NNLine!35,line width=0.3pt},
    tick align=outside,tick style={NNInk!55,line width=0.55pt},
    tick label style={font=\figurefont\footnotesize,text=NNInk},
    label style={font=\figurefont\footnotesize,text=NNInk},
    xlabel={$z$},ylabel={$\sigma(z)$},
    ylabel style={at={(axis cs:0,\pgfkeysvalueof{/pgfplots/ymax})},
      anchor=south,yshift=2pt},
    title style={font=\figurefont\footnotesize\bfseries,text=NNInk,
      at={(0,1.16)},anchor=south west},
    legend cell align=left,
    legend image post style={xscale=1.35},
    legend style={font=\figurefont\footnotesize,text=NNInk,draw=none,fill=none,
      at={(0.5,-0.18)},anchor=north,legend columns=1,
      inner sep=1pt,row sep=2pt,
      /tikz/every even column/.append style={column sep=3mm}},
    clip=true}}

\begin{tikzpicture}[font=\figurefont\footnotesize]
  \begin{axis}[ffn activation axis,at={(0,0)},xmin=-2,xmax=2,ymin=-2.3,ymax=2.3,xtick={-2,-1,1,2},ytick={-2,-1,1,2}]
    \addplot[NNHidden,ffn curve] table[x=z,y=identity] {figures/feedforward-activation-curves.dat};
    \addlegendentry{Identity}
  \end{axis}
  \begin{axis}[ffn activation axis,at={(85mm,0)},
    xmin=-2,xmax=2,ymin=-0.2,ymax=1.25,
    xtick={-2,-1,1,2},ytick={1}]
    \addplot[NNInput,ffn curve] coordinates {(-2,0) (0,0)};
    \addlegendentry{阶跃}
    \addplot[NNInput,ffn curve,forget plot] coordinates {(0,1) (2,1)};
    \addplot[only marks,mark=o,mark size=2pt,draw=NNInput,
      fill=white,forget plot] coordinates {(0,0)};
    \addplot[only marks,mark=*,mark size=2pt,draw=NNInput,
      fill=NNInput,forget plot] coordinates {(0,1)};
  \end{axis}
\end{tikzpicture}
\endgroup

  \caption{恒等函数与阶跃函数的比较。左图保留输入，斜率恒为$1$；右图按零阈值输出$0$或$1$，空心点和实心点区分$z=0$两侧的端点约定。横轴为预激活值$z$，纵轴为输出；阶跃图中的跳跃不能通过普通局部导数传递学习信号。}
  \label{fig:ffn-activation-basics}
\end{figure*}

\subsubsection{双侧饱和函数}
\label{sec:ffn-activation-saturating}

当一个单元需要产生概率、控制系数或有界状态时，把任意实数压到有限区间是直接的选择。\term{Logistic函数}产生$(0,1)$内的响应，实践中常直接称为Sigmoid；\term{双曲正切函数}（hyperbolic tangent, Tanh）产生$(-1,1)$内的响应。Sigmoid也可泛指具有S形曲线的一类函数，这里用$s$专门表示Logistic函数：
\begin{equation}
  \begin{aligned}
    s(z)&=\frac{1}{1+\exp(-z)},
    &s'(z)&=s(z)(1-s(z)),\\
    \tanh(z)&=\frac{\exp(z)-\exp(-z)}{\exp(z)+\exp(-z)},
    &\tanh'(z)&=1-\tanh^2(z)\text{。}
  \end{aligned}
  \label{eq:ffn-saturating-activations}
\end{equation}
当$|z|$很大时，输出接近两端的极限，导数趋近于零。这种\term{饱和}（saturation）意味着输入继续变化，输出却几乎不变。Sigmoid导数至多为$1/4$；Tanh在原点的导数为$1$，并满足$\tanh(z)=2s(2z)-1$。Tanh关于原点对称，可以产生正负响应，但只有对称输入分布等条件才能进一步推出输出均值为零。

Sigmoid适合二分类概率、多标签概率和取值在$(0,1)$的控制门；它不会自动赋予任意隐藏坐标概率意义。Tanh适合需要正负有界状态的模块，例如后续循环网络中的候选状态。两者用于许多层连续串联时，饱和区的小导数可能参与梯度衰减；传播中还包含权重矩阵，不能只凭激活导数判断整个梯度必然变小。

若希望避免指数运算，可以用分段线性函数近似上述有界响应。\term{硬Sigmoid}（hard sigmoid，也称Hard-Logistic）和\term{硬Tanh}（hard tanh）将线性段之外的值截到边界：
\begin{equation}
  \begin{aligned}
    \operatorname{HardSigmoid}(z)
      &=\min\{1,\max\{0,\tfrac14z+\tfrac12\}\},\\
    \operatorname{HardTanh}(z)&=\min\{1,\max\{-1,z\}\}\text{。}
  \end{aligned}
  \label{eq:ffn-hard-saturating-activations}
\end{equation}
本节的HardSigmoid在$(-2,2)$内斜率为$1/4$，区间外为$0$；HardTanh在$(-1,1)$内斜率为$1$，区间外为$0$，其负端输出是$-1$。分段边界处的导数需要约定。HardSigmoid的线性段并没有唯一标准，一些实现采用$z/6+1/2$并在$\pm3$处截断；移植模型时必须核对具体定义\scite{pytorchHardsigmoid}

两种硬近似可用于受计算资源限制的门控或范围约束，代价是拐点不可微，截断区内还会出现严格为零的梯度。\term{Softsign函数}提供另一种有界响应，用绝对值而非指数控制幅度：
\begin{equation}
  \operatorname{Softsign}(z)=\frac{z}{1+|z|},
  \qquad
  \operatorname{Softsign}'(z)=\frac{1}{(1+|z|)^2}\text{。}
  \label{eq:ffn-softsign}
\end{equation}
它的值域也是$(-1,1)$，一阶导数连续，尾部导数以多项式速度趋零。它可用于需要有界正负响应的隐藏层或状态映射；与Tanh相比，尾部衰减方式不同，并不意味着在所有输入尺度下梯度都更大。

\begin{figure*}[htb]
  \centering
  % 原创解析曲线。采样数据由 feedforward-curves.py 生成。
\begingroup
% 局部统一的科学绘图样式；曲线保持颜色和线型双重编码。
\pgfplotsset{
  ffn curve/.style={line width=1.3pt,line cap=round,line join=round},
  ffn reference curve/.style={NNInk!58,line width=0.95pt},
  ffn activation axis/.style={
    width=72mm,height=59mm,anchor=south west,
    axis lines=middle,
    axis line style={NNInk!65,line width=0.65pt},
    grid=major,grid style={NNLine!35,line width=0.3pt},
    tick align=outside,tick style={NNInk!55,line width=0.55pt},
    tick label style={font=\figurefont\footnotesize,text=NNInk},
    label style={font=\figurefont\footnotesize,text=NNInk},
    xlabel={$z$},ylabel={$\sigma(z)$},
    ylabel style={at={(axis cs:0,\pgfkeysvalueof{/pgfplots/ymax})},
      anchor=south,yshift=2pt},
    title style={font=\figurefont\footnotesize\bfseries,text=NNInk,
      at={(0,1.16)},anchor=south west},
    legend cell align=left,
    legend image post style={xscale=1.35},
    legend style={font=\figurefont\footnotesize,text=NNInk,draw=none,fill=none,
      at={(0.5,-0.18)},anchor=north,legend columns=1,
      inner sep=1pt,row sep=2pt,
      /tikz/every even column/.append style={column sep=3mm}},
    clip=true}}

\begin{tikzpicture}[font=\figurefont\footnotesize]
  \begin{axis}[ffn activation axis,at={(0,0)},xmin=-4,xmax=4,ymin=-1.25,ymax=1.2,xtick={-4,-2,2,4},ytick={-1,0.5,1}]
    \addplot[NNInput,ffn curve] table[x=z,y=sigmoid] {figures/feedforward-activation-curves.dat};
    \addlegendentry{Sigmoid}
    \addplot[NNHidden,ffn curve,dashed] table[x=z,y=tanh] {figures/feedforward-activation-curves.dat};
    \addlegendentry{Tanh}
    \addplot[NNOutput,ffn curve,densely dotted] table[x=z,y=softsign] {figures/feedforward-activation-curves.dat};
    \addlegendentry{Softsign}
  \end{axis}
  \begin{axis}[ffn activation axis,at={(85mm,0)},xmin=-4,xmax=4,ymin=-1.25,ymax=1.2,xtick={-4,-2,2,4},ytick={-1,0.5,1}]
    \addplot[NNInput,ffn curve] table[x=z,y=sigmoid] {figures/feedforward-activation-curves.dat};
    \addlegendentry{Sigmoid}
    \addplot[NNHidden,ffn curve,dashed] table[x=z,y=tanh] {figures/feedforward-activation-curves.dat};
    \addlegendentry{Tanh}
    \addplot[NNGradient,ffn curve,dashdotted] table[x=z,y=hardsigmoid] {figures/feedforward-activation-curves.dat};
    \addlegendentry{HardSigmoid}
    \addplot[NNOutput,ffn curve,densely dotted] table[x=z,y=hardtanh] {figures/feedforward-activation-curves.dat};
    \addlegendentry{HardTanh}
  \end{axis}
\end{tikzpicture}
\endgroup

  \caption{双侧饱和激活函数。左图比较Logistic／Sigmoid、Tanh与Softsign；右图将HardSigmoid和HardTanh与对应平滑函数比较。HardSigmoid采用式\eqref{eq:ffn-hard-saturating-activations}中的斜率$1/4$及截断点$\pm2$。各函数都限制输出范围，尾部趋近方式与梯度是否严格为零有所不同。}
  \label{fig:ffn-activations}
\end{figure*}

\subsubsection{整流函数及其变体}
\label{sec:ffn-activation-rectifiers}

ReLU采用前章的$\rho(z)=\max\{0,z\}$，正半轴斜率为$1$，负半轴输出和导数都为$0$。它只需比较与取值，能够产生严格为零的激活，可作为全连接或卷积隐藏层的基础选择。如果某单元对所有训练样本都长期落在负半轴，其输入权重可能难以得到数据梯度。这种失活要结合参数、数据和更新方式判断；在较深层，上游表示本身仍可能改变。

让负半轴保留非零斜率，可以避免该区域完全阻断局部梯度。\term{带泄漏的ReLU}（leaky ReLU）使用固定$0<\alpha<1$，\term{参数化ReLU}（parametric ReLU, PReLU）则把负半轴斜率作为待学习参数：
\begin{equation}
  \begin{aligned}
    \operatorname{LeakyReLU}(z)
      &=\max\{0,z\}+\alpha\min\{0,z\},\\
    \operatorname{PReLU}_j(z)
      &=\max\{0,z\}+\alpha_j\min\{0,z\}\text{。}
  \end{aligned}
  \label{eq:ffn-leaky-prelu}
\end{equation}
PReLU的$\alpha_j$可以按单元、通道设置，也可以在一层内共享；若不另加约束，它并不天然位于$(0,1)$。给定上游梯度，斜率参数的局部求导规则是$\partial\operatorname{PReLU}_j(z)/\partial\alpha_j=\min\{0,z\}$，共享斜率时要累加各使用位置的贡献。PReLU增加的参数量取决于共享方式，通常每个单元、通道或层只增加一个斜率参数\scite{he2015rectifiers}

Leaky ReLU适合希望保留负响应且不增加待学习参数的隐藏层，PReLU适合允许数据调整负响应强度的模块。\term{随机化带泄漏ReLU}（randomized leaky ReLU, RReLU）沿用同一分段形式，但训练时从预先规定的区间$[l,u]$均匀采样负斜率$\alpha$，其中$0<l<u<1$，预测时通常采用$(l+u)/2$。这种训练扰动可作为正则化设计的一部分，不能把训练阶段画成唯一固定曲线；它是否改善预测表现仍需评价\scite{xu2015rectified}

如果部署还要求限制正响应，可以使用\term{ReLU6函数}，即
\begin{equation}
  \operatorname{ReLU6}(z)=\min\{6,\max\{0,z\}\}\text{。}
  \label{eq:ffn-relu6}
\end{equation}
它在$(0,6)$内保留斜率$1$，在两侧导数为$0$，常与移动端模型及有限数值范围的设计联系在一起。MobileNetV2使用了这类有界整流响应，但其线性瓶颈层还刻意省略非线性，说明激活要按层的职责选择\scite{sandler2018mobilenetv2}ReLU6不是处处保持正半轴梯度的ReLU；超过上界的输入也会进入饱和区。

\begin{figure*}[htb]
  \centering
  % 原创解析曲线。采样数据由 feedforward-curves.py 生成。
\begingroup
% 局部统一的科学绘图样式；曲线保持颜色和线型双重编码。
\pgfplotsset{
  ffn curve/.style={line width=1.3pt,line cap=round,line join=round},
  ffn reference curve/.style={NNInk!58,line width=0.95pt},
  ffn activation axis/.style={
    width=72mm,height=59mm,anchor=south west,
    axis lines=middle,
    axis line style={NNInk!65,line width=0.65pt},
    grid=major,grid style={NNLine!35,line width=0.3pt},
    tick align=outside,tick style={NNInk!55,line width=0.55pt},
    tick label style={font=\figurefont\footnotesize,text=NNInk},
    label style={font=\figurefont\footnotesize,text=NNInk},
    xlabel={$z$},ylabel={$\sigma(z)$},
    ylabel style={at={(axis cs:0,\pgfkeysvalueof{/pgfplots/ymax})},
      anchor=south,yshift=2pt},
    title style={font=\figurefont\footnotesize\bfseries,text=NNInk,
      at={(0,1.16)},anchor=south west},
    legend cell align=left,
    legend image post style={xscale=1.35},
    legend style={font=\figurefont\footnotesize,text=NNInk,draw=none,fill=none,
      at={(0.5,-0.18)},anchor=north,legend columns=1,
      inner sep=1pt,row sep=2pt,
      /tikz/every even column/.append style={column sep=3mm}},
    clip=true}}

\begin{tikzpicture}[font=\figurefont\footnotesize]
  \begin{axis}[ffn activation axis,at={(0,0)},xmin=-4,xmax=4,ymin=-1.3,ymax=4.3,xtick={-4,-2,2,4},ytick={-1,2,4}]
    \addplot[NNInput,ffn curve] table[x=z,y=relu] {figures/feedforward-activation-curves.dat};
    \addlegendentry{ReLU}
    \addplot[NNHidden,ffn curve,dashed] table[x=z,y=leaky] {figures/feedforward-activation-curves.dat};
    \addlegendentry{Leaky: $0.1$}
    \addplot[NNOutput,ffn curve,dashdotted] table[x=z,y=prelu] {figures/feedforward-activation-curves.dat};
    \addlegendentry{PReLU: $0.25$}
    \addplot[NNGradient,ffn curve,densely dotted] table[x=z,y=rrelu] {figures/feedforward-activation-curves.dat};
    \addlegendentry{RReLU: $0.2$}
  \end{axis}
  \begin{axis}[ffn activation axis,at={(85mm,0)},xmin=-2,xmax=8,ymin=-0.4,ymax=8.4,xtick={-2,2,4,6,8},ytick={2,4,6,8}]
    \addplot[NNInput,ffn curve] table[x=z,y=relu] {figures/feedforward-activation-curves.dat};
    \addlegendentry{ReLU}
    \addplot[NNHidden,ffn curve,dashed] table[x=z,y=relu6] {figures/feedforward-activation-curves.dat};
    \addlegendentry{ReLU6}
  \end{axis}
\end{tikzpicture}
\endgroup

  \caption{整流函数与负斜率变体。左图固定Leaky ReLU的$\alpha=0.1$、PReLU的$\alpha_j=0.25$，并以$\alpha=0.2$画出RReLU训练阶段的一种可能曲线；三者的形状差异来自斜率选择，PReLU与RReLU的区别主要在参数取得方式。右图比较ReLU与ReLU6，展示$z>6$时的截断。两图横轴范围不同。}
  \label{fig:ffn-activation-rectifiers}
\end{figure*}

有界函数与整流函数的差异集中在表~\ref{tab:ffn-activation-bounded}中。Sigmoid、Tanh和Softsign限制两端幅度，ReLU只截断负半轴，ReLU6还限制正端；这种范围差异会直接改变哪些输入能够传递梯度。

\begin{table*}[htb]
  \centering
  \small
  \begin{tabularx}{\linewidth}{@{}l l X X@{}}
    \toprule
    函数 & 输出范围 & 导数与参数特点 & 适用场景及主要边界 \\
    \midrule
    Identity & $\mathbb{R}$ & 导数为$1$ & 回归输出、线性投影；不增加非线性 \\
    阶跃 & $\{0,1\}$ & 非阈值处导数为$0$ & 硬判定；不能直接靠普通反向传播训练 \\
    Sigmoid & $(0,1)$ & 两端导数趋零 & 二分类、多标签、控制门 \\
    Tanh & $(-1,1)$ & 两端导数趋零 & 有界正负状态；深层饱和需关注 \\
    HardSigmoid & $[0,1]$ & 截断区导数为$0$ & 简单门控；需核对斜率与截断点 \\
    HardTanh & $[-1,1]$ & 截断区导数为$0$ & 有界状态；拐点与严格零梯度 \\
    Softsign & $(-1,1)$ & 一阶连续，尾部导数趋零 & 有界正负表示；尾部衰减方式不同 \\
    ReLU & $[0,+\infty)$ & 负半轴为$0$，正半轴为$1$ & 全连接、卷积隐藏层；可能失活 \\
    Leaky ReLU & $\mathbb{R}$ & 固定正负斜率 & 保留负梯度；斜率需设定 \\
    PReLU & 依赖$\alpha_j$ & 负斜率参与学习 & 学习响应形状；需说明参数共享方式 \\
    RReLU & $\mathbb{R}$ & 训练时随机负斜率 & 隐藏层训练扰动；预测规则不同 \\
    ReLU6 & $[0,6]$ & 两侧截断，中间斜率为$1$ & 有限范围、移动端；正端也会饱和 \\
    \bottomrule
  \end{tabularx}
  \caption{有界响应与整流激活的比较。Leaky ReLU和RReLU的范围按正负斜率均为正且非零书写；PReLU若不限制斜率，其输出范围随参数改变。所有分段函数都需在不可微边界约定求导规则。}
  \label{tab:ffn-activation-bounded}
\end{table*}

即使都用于限制输出幅度，求导行为也可能不同。Tanh和HardTanh都提供有界正负响应；按解析定义，前者在有限输入下保留非零导数，后者在截断区的导数严格为零。选择范围约束时，还需判断训练中的预激活值是否经常落到这些区域。

\subsubsection{指数型与平滑整流函数}
\label{sec:ffn-activation-exponential}

固定负斜率让负输出无限下降，另一种做法是使负响应逐渐接近一个有限下界。\term{指数线性单元}（exponential linear unit, ELU）在正半轴保留输入，在负半轴使用指数曲线：
\begin{equation}
  \operatorname{ELU}_{\alpha}(z)=
  \begin{cases}
    z,&z>0,\\
    \alpha(\exp(z)-1),&z\leq0,
  \end{cases}
  \qquad \alpha>0\text{。}
  \label{eq:ffn-elu}
\end{equation}
负半轴导数为$\alpha\exp(z)$，对有限负输入非零，但在$z\to-\infty$时趋零；输出趋近$-\alpha$。ELU连续，只有$\alpha=1$时左右一阶导数才在原点一致，而且此时也不具有所有阶的连续导数。ELU可用于希望保留负值并控制负端幅度的隐藏层；输出存在负值并不保证任意数据下的均值恰好为零\scite{clevert2016elu}

\term{缩放指数线性单元}（scaled exponential linear unit, SELU）进一步把ELU整体乘以固定系数：
\begin{equation}
  \operatorname{SELU}(z)=\kappa
  \begin{cases}
    z,&z>0,\\
    \alpha(\exp(z)-1),&z\leq0,
  \end{cases}
  \label{eq:ffn-selu}
\end{equation}
其中常用常数为$\alpha\approx1.6733$、$\kappa\approx1.0507$。它用于\term{自归一化网络}（self-normalizing neural network），即在规定假设下让层间激活的均值与方差趋向受控范围的前馈网络。相应分析依赖输入与权重的统计条件；常配合输入标准化、方差$1/\mathrm{fan\text{-}in}$的初始化，以及为保持统计量设计的AlphaDropout。它不是把ReLU换成SELU后就无条件得到的归一化保证，也不能随意混用初始化和丢弃规则\scite{klambauer2017selu}

如果目标是输出严格正值而又保持平滑，可以使用\term{Softplus函数}：
\begin{equation}
  \operatorname{Softplus}(z)=\log(1+\exp(z)),
  \qquad \operatorname{Softplus}'(z)=s(z)\text{。}
  \label{eq:ffn-softplus}
\end{equation}
它在正端接近$z$，在负端趋近$0$，但有限输入下输出严格大于$0$，不会产生ReLU式的精确零激活。Softplus可作为平滑隐藏激活，也适合将任意得分映射为正尺度、强度或方差参数；需要严格正下界时还可加上一个固定小常数。计算时使用$\max\{z,0\}+\log(1+\exp(-|z|))$可避免指数溢出。因此是否使用Softplus应按平滑性和正值约束判断，不能仅凭某个分类实验就认定它在所有场景中都不合适。

\begin{figure}[htbp]
  \centering
  % 原创解析曲线。采样数据由 feedforward-curves.py 生成。
\begingroup
% 局部统一的科学绘图样式；曲线保持颜色和线型双重编码。
\pgfplotsset{
  ffn curve/.style={line width=1.3pt,line cap=round,line join=round},
  ffn reference curve/.style={NNInk!58,line width=0.95pt},
  ffn activation axis/.style={
    width=72mm,height=59mm,anchor=south west,
    axis lines=middle,
    axis line style={NNInk!65,line width=0.65pt},
    grid=major,grid style={NNLine!35,line width=0.3pt},
    tick align=outside,tick style={NNInk!55,line width=0.55pt},
    tick label style={font=\figurefont\footnotesize,text=NNInk},
    label style={font=\figurefont\footnotesize,text=NNInk},
    xlabel={$z$},ylabel={$\sigma(z)$},
    ylabel style={at={(axis cs:0,\pgfkeysvalueof{/pgfplots/ymax})},
      anchor=south,yshift=2pt},
    title style={font=\figurefont\footnotesize\bfseries,text=NNInk,
      at={(0,1.16)},anchor=south west},
    legend cell align=left,
    legend image post style={xscale=1.35},
    legend style={font=\figurefont\footnotesize,text=NNInk,draw=none,fill=none,
      at={(0.5,-0.18)},anchor=north,legend columns=1,
      inner sep=1pt,row sep=2pt,
      /tikz/every even column/.append style={column sep=3mm}},
    clip=true}}

\begin{tikzpicture}[font=\figurefont\footnotesize]
  \begin{axis}[ffn activation axis,width=112mm,at={(0,0)},xmin=-4,xmax=4,ymin=-2.1,ymax=4.6,xtick={-4,-2,2,4},ytick={-2,2,4}]
    \addplot[ffn reference curve,densely dotted] table[x=z,y=relu] {figures/feedforward-activation-curves.dat};
    \addlegendentry{ReLU}
    \addplot[NNInput,ffn curve] table[x=z,y=elu] {figures/feedforward-activation-curves.dat};
    \addlegendentry{ELU}
    \addplot[NNHidden,ffn curve,dashed] table[x=z,y=selu] {figures/feedforward-activation-curves.dat};
    \addlegendentry{SELU}
    \addplot[NNOutput,ffn curve,dashdotted] table[x=z,y=softplus] {figures/feedforward-activation-curves.dat};
    \addlegendentry{Softplus}
  \end{axis}
\end{tikzpicture}
\endgroup

  \caption{ELU、SELU与Softplus的响应。ELU使用$\alpha=1$，SELU使用式\eqref{eq:ffn-selu}中的常数，ReLU作为共同参照。ELU与SELU允许负输出并在负端趋向有限下界；Softplus始终为正，且平滑地趋近正半轴上的线性响应。}
  \label{fig:ffn-activation-exponential}
\end{figure}

\subsubsection{自门控函数}
\label{sec:ffn-activation-self-gated}

ReLU根据输入的正负作硬切换，也可以让输入乘一个连续变化的系数。\term{Swish函数}使用输入自身计算保留系数：
\begin{equation}
  \operatorname{Swish}_{\beta}(z)=z s(\beta z),
  \qquad \beta\geq0\text{。}
  \label{eq:ffn-swish}
\end{equation}
$\beta$可以固定，也可以参与学习。$\beta=0$时函数为$z/2$；$\beta=1$时得到\term{Sigmoid加权线性单元}（sigmoid-weighted linear unit, SiLU）；当$\beta\to+\infty$时，对每个固定输入趋向ReLU。这里的变化描述响应曲线的极限，不能据此推出不同$\beta$的训练效果也连续地在两种模型之间变化\scite{ramachandran2017swish}

Swish的局部导数为
\begin{equation}
  \begin{aligned}
    \operatorname{Swish}_{\beta}'(z)
      &=s(\beta z)+\beta z s(\beta z)(1-s(\beta z)),\\
    \frac{\partial\operatorname{Swish}_{\beta}(z)}{\partial\beta}
      &=z^2s(\beta z)(1-s(\beta z))\text{。}
  \end{aligned}
  \label{eq:ffn-swish-derivatives}
\end{equation}
SiLU因而不是与Swish完全不同的函数，而是其固定参数实例。SiLU在强化学习函数逼近等场景中也有使用\scite{elfwing2018silu}SiLU可作为全连接、卷积或门控前馈模块中的平滑激活；在有限负区间内导数可能为负，负尾部导数又趋零。若学习$\beta$并要求保持非负，还需明确约束方式。

\begin{figure}[htbp]
  \centering
  % 原创解析曲线。采样数据由 feedforward-curves.py 生成。
\begingroup
% 局部统一的科学绘图样式；曲线保持颜色和线型双重编码。
\pgfplotsset{
  ffn curve/.style={line width=1.3pt,line cap=round,line join=round},
  ffn reference curve/.style={NNInk!58,line width=0.95pt},
  ffn activation axis/.style={
    width=72mm,height=59mm,anchor=south west,
    axis lines=middle,
    axis line style={NNInk!65,line width=0.65pt},
    grid=major,grid style={NNLine!35,line width=0.3pt},
    tick align=outside,tick style={NNInk!55,line width=0.55pt},
    tick label style={font=\figurefont\footnotesize,text=NNInk},
    label style={font=\figurefont\footnotesize,text=NNInk},
    xlabel={$z$},ylabel={$\sigma(z)$},
    ylabel style={at={(axis cs:0,\pgfkeysvalueof{/pgfplots/ymax})},
      anchor=south,yshift=2pt},
    title style={font=\figurefont\footnotesize\bfseries,text=NNInk,
      at={(0,1.16)},anchor=south west},
    legend cell align=left,
    legend image post style={xscale=1.35},
    legend style={font=\figurefont\footnotesize,text=NNInk,draw=none,fill=none,
      at={(0.5,-0.18)},anchor=north,legend columns=1,
      inner sep=1pt,row sep=2pt,
      /tikz/every even column/.append style={column sep=3mm}},
    clip=true}}

\begin{tikzpicture}[font=\figurefont\footnotesize]
  \begin{axis}[ffn activation axis,width=112mm,at={(0,0)},xmin=-4,xmax=4,ymin=-2.3,ymax=4.3,xtick={-4,-2,2,4},ytick={-2,2,4}]
    \addplot[NNInk,ffn curve,densely dotted] table[x=z,y=swish0] {figures/feedforward-activation-curves.dat};
    \addlegendentry{$\beta=0$}
    \addplot[NNHidden,ffn curve,dashed] table[x=z,y=swish05] {figures/feedforward-activation-curves.dat};
    \addlegendentry{$\beta=0.5$}
    \addplot[NNInput,ffn curve] table[x=z,y=silu] {figures/feedforward-activation-curves.dat};
    \addlegendentry{$\beta=1$ (SiLU)}
    \addplot[NNOutput,ffn curve,dashdotted] table[x=z,y=swish5] {figures/feedforward-activation-curves.dat};
    \addlegendentry{$\beta=5$}
    \addplot[ffn reference curve,dash pattern=on 5pt off 1pt on 1pt off 1pt] table[x=z,y=relu] {figures/feedforward-activation-curves.dat};
    \addlegendentry{ReLU}
  \end{axis}
\end{tikzpicture}
\endgroup

  \caption{Swish在不同$\beta$下的函数图像。$\beta=0$给出线性函数$z/2$，$\beta=1$就是SiLU，较大的$\beta$逐渐形成接近ReLU的响应。图中参数固定为$0,0.5,1,5$，ReLU用作极限参照；它不展示可学习$\beta$在某次训练中的实际轨迹。}
  \label{fig:ffn-activation-swish}
\end{figure}

\term{高斯误差线性单元}（Gaussian error linear unit, GELU）采用标准正态累积分布函数$\Phi$作为保留系数：
\begin{equation}
  \operatorname{GELU}(z)=z\Phi(z),
  \qquad \operatorname{GELU}'(z)=\Phi(z)+z\varphi(z)\text{。}
  \label{eq:ffn-smooth-gates}
\end{equation}
其中$\varphi(z)=\exp(-z^2/2)/\sqrt{2\pi}$是标准正态密度。GELU在全实数上平滑，但公式中的分布函数只定义一个确定变换，并不表示执行时随机丢弃输入。它常作为Transformer前馈模块的激活选择，平滑性本身不能保证找到更好的最优解。常见近似为$z s(1.702z)$，或者
\begin{equation}
  \operatorname{GELU}(z)\approx
  \frac z2\left(1+\tanh\!\left[\sqrt{\frac2\pi}
  (z+0.044715z^3)\right]\right)\text{。}
  \label{eq:ffn-gelu-approximation}
\end{equation}
使用哪种近似会改变实际数值与导数，复现模型时应记录具体形式。上述近似与精确定义的关系是数值近似，不是把GELU与SiLU视为同一个函数\scite{hendrycks2016gelu}

\term{Mish函数}进一步用Softplus和Tanh形成平滑的保留系数\scite{misra2019mish}其定义为
\begin{equation}
  \operatorname{Mish}(z)=z\tanh(\operatorname{Softplus}(z))\text{。}
  \label{eq:ffn-mish}
\end{equation}
令$t=\tanh(\operatorname{Softplus}(z))$，乘积与链式法则给出导数$t+zs(z)(1-t^2)$。Mish在负半轴保留一段负响应，负尾部趋近零，正端接近线性，可作为视觉网络或普通隐藏层的另一种平滑选择；原论文的比较依赖具体模型与数据，不能概括为普遍优于其他激活。

这些平滑函数需要指数、Tanh或正态分布函数计算。若希望用简单分段运算近似SiLU，可使用\term{硬Swish}（hard swish, Hardswish）。MobileNetV3在面向移动端的结构中使用这种响应\scite{howard2019mobilenetv3}其定义为
\begin{equation}
  \operatorname{Hardswish}(z)=z\frac{\operatorname{ReLU6}(z+3)}6
  =\begin{cases}
    0,&z\leq-3,\\
    z(z+3)/6,&-3<z<3,\\
    z,&z\geq3\text{。}
  \end{cases}
  \label{eq:ffn-hardswish}
\end{equation}
它在中间区间的导数为$(2z+3)/6$，在两端分别为$0$和$1$，在$\pm3$处不具有唯一导数。这里的门采用斜率$1/6$的硬Sigmoid形式，不能直接替换成本节前面斜率$1/4$的版本。实际速度取决于算子与硬件实现，而不只取决于公式中出现几个基本运算。

\begin{figure*}[htb]
  \centering
  % 原创解析曲线。采样数据由 feedforward-curves.py 生成。
\begingroup
% 局部统一的科学绘图样式；曲线保持颜色和线型双重编码。
\pgfplotsset{
  ffn curve/.style={line width=1.3pt,line cap=round,line join=round},
  ffn reference curve/.style={NNInk!58,line width=0.95pt},
  ffn activation axis/.style={
    width=72mm,height=59mm,anchor=south west,
    axis lines=middle,
    axis line style={NNInk!65,line width=0.65pt},
    grid=major,grid style={NNLine!35,line width=0.3pt},
    tick align=outside,tick style={NNInk!55,line width=0.55pt},
    tick label style={font=\figurefont\footnotesize,text=NNInk},
    label style={font=\figurefont\footnotesize,text=NNInk},
    xlabel={$z$},ylabel={$\sigma(z)$},
    ylabel style={at={(axis cs:0,\pgfkeysvalueof{/pgfplots/ymax})},
      anchor=south,yshift=2pt},
    title style={font=\figurefont\footnotesize\bfseries,text=NNInk,
      at={(0,1.16)},anchor=south west},
    legend cell align=left,
    legend image post style={xscale=1.35},
    legend style={font=\figurefont\footnotesize,text=NNInk,draw=none,fill=none,
      at={(0.5,-0.18)},anchor=north,legend columns=1,
      inner sep=1pt,row sep=2pt,
      /tikz/every even column/.append style={column sep=3mm}},
    clip=true}}

\begin{tikzpicture}[font=\figurefont\footnotesize]
  \begin{axis}[ffn activation axis,at={(0,0)},xmin=-4,xmax=4,ymin=-0.6,ymax=4.3,xtick={-4,-2,2,4},ytick={2,4}]
    \addplot[NNInput,ffn curve] table[x=z,y=silu] {figures/feedforward-activation-curves.dat};
    \addlegendentry{SiLU}
    \addplot[NNHidden,ffn curve,dashed] table[x=z,y=gelu] {figures/feedforward-activation-curves.dat};
    \addlegendentry{GELU}
    \addplot[NNOutput,ffn curve,dashdotted] table[x=z,y=mish] {figures/feedforward-activation-curves.dat};
    \addlegendentry{Mish}
    \addplot[NNGradient,ffn curve,densely dotted] table[x=z,y=hardswish] {figures/feedforward-activation-curves.dat};
    \addlegendentry{Hardswish}
  \end{axis}
  \begin{axis}[ffn activation axis,at={(85mm,0)},xmin=-3.5,xmax=1,ymin=-0.45,ymax=0.95,xtick={-3,-2,-1,1},ytick={-0.3,0.3,0.6,0.9}]
    \addplot[NNInput,ffn curve] table[x=z,y=silu] {figures/feedforward-activation-curves.dat};
    \addlegendentry{SiLU}
    \addplot[NNHidden,ffn curve,dashed] table[x=z,y=gelu] {figures/feedforward-activation-curves.dat};
    \addlegendentry{GELU}
    \addplot[NNOutput,ffn curve,dashdotted] table[x=z,y=mish] {figures/feedforward-activation-curves.dat};
    \addlegendentry{Mish}
    \addplot[NNGradient,ffn curve,densely dotted] table[x=z,y=hardswish] {figures/feedforward-activation-curves.dat};
    \addlegendentry{Hardswish}
  \end{axis}
\end{tikzpicture}
\endgroup

  \caption{SiLU、GELU、Mish与Hardswish。左图展示整体响应，右图放大原点附近的负值和过渡区域；GELU按标准正态累积分布函数采样。它们在正端接近线性，但负区间的形状和拐点不同。SiLU、GELU与Mish平滑，Hardswish含有分段边界；两图使用不同坐标范围以显示细节。}
  \label{fig:ffn-activation-smooth}
\end{figure*}

\subsubsection{Maxout与Softmax}
\label{sec:ffn-activation-vector}

以上函数分别作用于一个标量。若希望一个单元从若干候选响应中选择最大者，可以使用\term{Maxout单元}（maxout unit）：它先计算$K$个仿射响应，再取最大值，响应的分段形状由权重学习得到。对于输入$\bm{u}\in\mathbb{R}^{m}$，其定义为
\begin{equation}
  z_k=\bm{w}_k^{\top}\bm{u}+b_k,
  \qquad h=\max_{1\leq k\leq K}z_k\text{，}
  \label{eq:ffn-maxout}
\end{equation}
其中每个$\bm{w}_k\in\mathbb{R}^{m}$、$b_k\in\mathbb{R}$都参与学习。只有一个严格最大响应时，梯度传给这一响应；多个响应并列最大时需约定分配方式。固定一个响应为零、另一个为$\bm{w}^{\top}\bm{u}+b$，就得到ReLU这一特殊形式。一个Maxout单元对输入表示为凸分段仿射函数，不意味着包含多个Maxout层的完整网络只能表示凸函数\scite{goodfellow2013maxout}

Maxout适合希望学习分段响应形状的隐藏单元，也曾与Dropout联合研究。相对于一个普通仿射响应，它有$K$套权重与偏置，参数和响应计算量通常约增至$K$倍，而不是无论$K$多大都只增加一倍。它不包含固定的负半轴零梯度区域，但某些分支仍可能长期不是最大者而收不到梯度。它还不再是式\eqref{eq:ffn-hidden-layers}中逐坐标作用于一个仿射响应的激活，反向计算需使用最大值选择对应的规则。

若目的不是选择一个最大数值，而是把$K$个得分转换为互斥事件的概率，可以使用\term{Softmax变换}。对$\bm{z}\in\mathbb{R}^{K}$，其第$k$个输出为
\begin{equation}
  p_k=\operatorname{softmax}(\bm{z})_k
  =\frac{\exp(z_k)}{\sum_{j=1}^{K}\exp(z_j)},
  \qquad \sum_{k=1}^{K}p_k=1\text{。}
  \label{eq:ffn-softmax-activation}
\end{equation}
每个概率依赖全部输入坐标，所以不存在脱离其他得分的唯一一元Softmax曲线。其雅可比为$\partial p_k/\partial z_j=p_k(\mathbb{I}\{k=j\}-p_j)$。Softmax适合互斥多分类输出和需要归一化权重的模块；多标签输出则通常分别使用Sigmoid。训练时把Softmax与交叉熵合并，可以得到更直接且数值稳定的梯度，下一小节将展开。

\begin{figure*}[htb]
  \centering
  % 原创解析曲线。采样数据由 feedforward-curves.py 生成。
\begingroup
% 局部统一的科学绘图样式；曲线保持颜色和线型双重编码。
\pgfplotsset{
  ffn curve/.style={line width=1.3pt,line cap=round,line join=round},
  ffn reference curve/.style={NNInk!58,line width=0.95pt},
  ffn activation axis/.style={
    width=72mm,height=59mm,anchor=south west,
    axis lines=middle,
    axis line style={NNInk!65,line width=0.65pt},
    grid=major,grid style={NNLine!35,line width=0.3pt},
    tick align=outside,tick style={NNInk!55,line width=0.55pt},
    tick label style={font=\figurefont\footnotesize,text=NNInk},
    label style={font=\figurefont\footnotesize,text=NNInk},
    xlabel={$z$},ylabel={$\sigma(z)$},
    ylabel style={at={(axis cs:0,\pgfkeysvalueof{/pgfplots/ymax})},
      anchor=south,yshift=2pt},
    title style={font=\figurefont\footnotesize\bfseries,text=NNInk,
      at={(0,1.16)},anchor=south west},
    legend cell align=left,
    legend image post style={xscale=1.35},
    legend style={font=\figurefont\footnotesize,text=NNInk,draw=none,fill=none,
      at={(0.5,-0.18)},anchor=north,legend columns=1,
      inner sep=1pt,row sep=2pt,
      /tikz/every even column/.append style={column sep=3mm}},
    clip=true}}

\begin{tikzpicture}[font=\figurefont\footnotesize]
  \begin{axis}[ffn activation axis,at={(0,0)},xmin=-3,xmax=3,ymin=-2,ymax=5.7,xtick={-2,-1,1,2},ytick={-1,1,3,5},ylabel={$h$}]
    \addplot[ffn reference curve,dashed] table[x=z,y=branch1] {figures/feedforward-activation-curves.dat};
    \addlegendentry{$-z-1$}
    \addplot[NNInput,line width=0.95pt,densely dotted] table[x=z,y=branch2] {figures/feedforward-activation-curves.dat};
    \addlegendentry{$z/2$}
    \addplot[NNGradient,line width=0.95pt,dashdotted] table[x=z,y=branch3] {figures/feedforward-activation-curves.dat};
    \addlegendentry{$2z-1$}
    \addplot[NNHidden,line width=1.55pt,line cap=round,line join=round] table[x=z,y=maxout] {figures/feedforward-activation-curves.dat};
    \addlegendentry{Maxout}
  \end{axis}
  \begin{axis}[ffn activation axis,at={(85mm,0)},xmin=-4,xmax=4,ymin=0,ymax=1.12,xtick={-4,-2,2,4},ytick={0.5,1},ylabel={$p_k$}]
    \addplot[NNHidden,ffn curve] table[x=z,y=probability1] {figures/feedforward-activation-curves.dat};
    \addlegendentry{$p_1$}
    \addplot[NNInput,ffn curve,dashed] table[x=z,y=probability23] {figures/feedforward-activation-curves.dat};
    \addlegendentry{$p_2=p_3$}
  \end{axis}
\end{tikzpicture}
\endgroup

  \caption{Maxout和Softmax的一维观察方式。左图取三个仿射候选$-z-1$、$0.5z$和$2z-1$，粗实线是它们的最大值包络；它展示固定参数后沿一维输入的Maxout响应。右图固定三类得分为$(z,0,0)$，绘出三项Softmax概率，其中$p_2=p_3$，用同一条曲线标注。这些切片不能代表所有输入与参数下的唯一函数图像。}
  \label{fig:ffn-activation-vector}
\end{figure*}

\subsubsection{激活函数的比较与选择}
\label{sec:ffn-activation-selection}

表~\ref{tab:ffn-activation-bounded}归纳了前面的有界响应与整流函数，表~\ref{tab:ffn-activation-smooth}补充指数型、自门控与向量变换。比较时，先确定该层需要输出什么，再检查梯度和计算代价。表中的使用场景说明功能匹配；可学习参数和随机机制也需要与固定曲线区分。



\begin{table*}[htb]
  \centering
  \small
  \begin{tabularx}{\linewidth}{@{}l l X X@{}}
    \toprule
    函数 & 输出范围 & 导数与参数特点 & 适用场景及主要边界 \\
    \midrule
    ELU & $(-\alpha,+\infty)$ & 负端导数趋零，$\alpha$固定 & 允许负响应；零均值没有无条件保证 \\
    SELU & $(-\kappa\alpha,+\infty)$ & 固定缩放与负指数响应 & 自归一化MLP；需配套统计条件 \\
    Softplus & $(0,+\infty)$ & 平滑，导数为Sigmoid & 正尺度输出、平滑隐藏层；无精确零 \\
    Swish／SiLU & 依赖$\beta$ & $\beta$固定或学习，SiLU取$1$ & 平滑隐藏层、前馈模块；局部可非单调 \\
    GELU & 有负下界，正端无界 & 平滑，精确式或近似式 & Transformer前馈模块；非随机采样 \\
    Mish & 有负下界，正端无界 & 平滑，复合运算 & 视觉或其他隐藏层；收益需验证 \\
    Hardswish & $[-3/8,+\infty)$ & 两个分段边界 & 移动端响应；速度取决于实现 \\
    Maxout & 依赖候选仿射响应 & 选中分支获梯度，$K$套参数 & 学习分段函数；分支可能少获梯度 \\
    Softmax & 概率单纯形内部 & 联合依赖全部得分 & 互斥类别、归一化权重；非逐坐标激活 \\
    \bottomrule
  \end{tabularx}
  \caption{指数型、自门控与向量变换的比较。Swish在$\beta=0$时是线性函数；其余正$\beta$情形具有有限负下界。Softmax按$K\geq2$书写，输出各项严格为正且总和为$1$。表中的平滑性描述固定参数的标量函数，不能替代整条计算路径的可微性检查。}
  \label{tab:ffn-activation-smooth}
\end{table*}

若该层预测概率或正尺度，输出约束通常决定选择：互斥类别用Softmax，多个可共存标签分别用Sigmoid，正尺度可以用Softplus。若该层构造隐藏表示，则需要比较非线性、活动范围与梯度。ReLU提供简单的分段响应，Leaky ReLU与PReLU保留负梯度，GELU与SiLU提供平滑过渡；SELU则是与特定初始化和统计条件配套的设计，不能只当作另一条任意替换的曲线。

计算成本也必须放回实现环境。分段线性函数少用指数等运算，但向量化、融合算子和内存访问可能改变实际耗时。平滑函数可能便于需要连续导数的模型，但不会自动避免小梯度或改善泛化；增加可学习斜率也可能改变初始化时的合适尺度。比较预测性能时，应固定数据划分、结构和评价口径，并说明是否重新调节学习率与初始化，才能区分激活变化和配套设置变化的影响。

% 历史开篇的两条完整边注文献需要与段落一起容纳。
\Needspace{16\baselineskip}
\subsubsection{从阈值单元到平滑门控：发展脉络}
\label{sec:ffn-activation-history}

McCulloch与Pitts在1943年把神经活动的“全或无”性质写成阈值单元；Rosenblatt在1958年的感知机工作中进一步讨论了连接如何从样本中调整\scite{mcculloch1943,rosenblatt1958}这些早期模型用单元的开关组合表达判定。硬阈值便于表达逻辑关系，但其普通导数几乎处处为零，无法直接承担本章所用的多层梯度学习。

当隐藏层需要通过损失学习表示时，连续响应提供了可用的局部导数。Rumelhart、Hinton与Williams在1986年的反向传播论文中使用Logistic响应，展示了隐藏单元怎样随训练形成特征。这一工作将可微的连续响应与隐藏表示的学习结合起来，使激活函数的选择与训练机制紧密相连\scite{rumelhart1986}

Sigmoid与Tanh随后也暴露出深层训练的尺度问题：两端饱和时，导数会变小，Sigmoid在原点附近又有非零的输出基线。2010年Glorot与Bengio的工作把激活响应与初始化放在一起分析，说明不能只要求单个函数可导，还要检查数值和梯度经过多层之后是否仍有合适尺度。Tanh把输出中心移到零附近，但仍保留两端饱和的边界\scite{glorot2010initialization}

这项工作的对照还表明，即使保留同一种激活，改变权重初始化方式也会改变各层活动值与梯度的尺度。选择激活函数时，因而还要把初始化和层深放在一起考察：一条曲线的局部斜率不足以描述整条网络的训练状态。第~\ref{sec:ffn-initialization}节将从信号的二阶矩出发解释这种配合，区分Tanh的中心响应与ReLU的半轴截断。

整流函数让设计重点转向稀疏响应和正半轴上的线性传播。Nair与Hinton在2010年研究了受限玻尔兹曼机中的整流单元；Glorot、Bordes与Bengio在2011年进一步比较深层前馈网络的整流与Tanh响应。这两项工作分别在生成模型与深层判别网络中展示了整流响应的作用。正半轴不饱和缓解了一类梯度衰减，负半轴的零导数又促使后续方法保留泄漏斜率或允许负响应\scite{nair2010rectifiers,glorot2011rectifiers}

较新的分支把注意力转向平滑过渡与输入相关的门控。2016年的GELU按输入幅度平滑加权；2017年的Swish工作通过函数搜索研究自门控响应，并与ReLU比较。两者都保留较大的正响应，同时用平滑过渡替代正负区域之间的硬切换\scite{hendrycks2016gelu,ramachandran2017swish}

本节介绍的带斜率变体、指数响应和硬近似，分别延伸了可学习形状、信号尺度与计算成本这些方向。历史上的改进解决的是特定结构和训练设置中的问题；今天选择隐藏激活时仍需核对这些条件。Sigmoid和Softmax在概率输出中的用途，则由预测对象决定，并不因隐藏层采用ReLU或平滑门控而消失。

\subsection{输出和学习目标}
\label{sec:ffn-output-objectives}

隐藏层得到的是一组可供后续使用的数值，输出层则要回答任务中的具体问题：预测一个连续数值、判断一个事件是否成立、从多个类别中选一个，还是生成供其他模块使用的表示？这些问题对应不同的输出空间。选定输出变换$g$之后，还需用损失函数说明什么预测算好、错误怎样计分。下面始终把最后一层的原始得分记为$\bm{a}$，区分原始得分、预测结果与训练损失。

图~\ref{fig:ffn-output-heads}把五种任务放在同一个前馈结构之后比较：左侧逐层形成隐藏表示，右侧按任务选择一个输出头，即把隐藏表示转换成所需预测的最后一组计算。各头有各自的权重、偏置和训练目标；图中的分支表示可替换的方案，并不表示必须同时训练五个任务。

\begin{figure*}[!htb]
  \centering
  % 按各输出头的内容高度留出相近的任务间隔。
  % 一个显式神经元骨干，加上五种可替换的任务输出头。
% 标量任务中心间隔 27 mm，向量任务间隔 32 mm，使内容之间均留约 8 mm。
% 分支起点是整组隐藏坐标；汇合线表示向量传递，不是求和算子。
\begin{tikzpicture}[x=1mm,y=1mm,font=\figurefont\footnotesize,
  head unit/.style={nn output,minimum size=5.5mm},
  head link/.style={nn flow,draw=NNOutput,line width=1.1pt},
  result/.style={anchor=west,align=left,text=NNInk,inner sep=0pt},
  task title/.style={nn heading,anchor=west,inner sep=0pt}]
  \node[nn heading,anchor=west] at (1,86) {共同的前馈网络};
  \node[nn heading,anchor=west] at (77,86) {按任务选择一个输出头};
  \draw[NNLine,line width=.65pt] (1,81)--(59,81);
  \draw[NNLine,line width=.65pt] (77,81)--(168,81);
  \node[nn annotation,align=left,anchor=north west] at (1,71)
    {输入逐层形成隐藏表示\\输出头决定预测的含义};

  % 3--4--4 的共同骨干与第三个输出头居中对齐。
  \begin{scope}[shift={(0,13)}]
    \foreach \i/\y in {1/7,2/-5,3/-17} \coordinate (x\i) at (6,\y);
    \foreach \i/\y in {1/13,2/1,3/-11,4/-23}{
      \coordinate (h\i) at (27,\y);
      \coordinate (k\i) at (48,\y);
    }
    \foreach \i in {1,2,3} \foreach \j in {1,2,3,4}
      \draw[nn link,draw=NNInput!58,shorten <=3.3mm,shorten >=3.3mm] (x\i)--(h\j);
    \foreach \i in {1,2,3,4} \foreach \j in {1,2,3,4}
      \draw[nn link,draw=NNHidden!58,shorten <=3.3mm,shorten >=3.3mm] (h\i)--(k\j);
    \foreach \i in {1,2,3}
      \node[nn input,minimum size=6mm] at (x\i) {};
    \foreach \i in {1,2,3,4}{
      \node[nn hidden,minimum size=6mm] at (h\i) {};
      \node[nn hidden,minimum size=6mm] at (k\i) {};
    }
    \node[nn annotation] at (6,24) {$\bm{x}$};
    \node[nn annotation] at (27,24) {$\bm{h}^{(1)}$};
    \node[nn annotation] at (48,24) {$\bm{h}^{(L)}$};
    \draw[NNHidden!60,line width=.85pt] (54,17)--(56,17)--(56,-27)--(54,-27);
    \draw[nn flow,draw=NNHidden] (56,-5)--(65,-5);
    \node[nn annotation,align=left,anchor=north west] at (1,-38)
      {圆点：一个神经元\\细线：可学习的连接};
    \node[nn annotation,align=left,anchor=north west,text width=57mm] at (1,-56)
      {各头读取整组$\bm{h}^{(L)}$，\\并各自学习输出权重和偏置。};
  \end{scope}
  \draw[NNLine,line width=.85pt] (65,62)--(65,-56);
  \foreach \y/\tag/\title in {62/a/回归,35/b/二分类,8/c/互斥多分类,-24/d/多标签分类,-56/e/表示输出}{
    \draw[head link,draw=NNHidden!75] (65,\y)--(82,\y);
    \node[task title] at (78,\y+13) {(\tag)\ \title};
  }

  % 标量输出头。
  \foreach \y/\tag/\transform/\prediction/\loss in {
    62/a/恒等/\widehat y=a\in\mathbb{R}/平方损失,
    35/b/Sigmoid/p=s(a)/二元交叉熵}{
    \begin{scope}[shift={(0,\y)}]
      \node[head unit] (score\tag) at (86,0) {$a$};
      \node[nn transform,minimum width=22mm] (transform\tag) at (114,0) {\transform};
      \draw[head link] (score\tag)--(transform\tag);
      \draw[head link] (transform\tag.east)--(132,0);
      \node[result] at (136,3) {$\prediction$};
      \node[result,text=BookMuted] at (136,-3) {\loss};
    \end{scope}
  }

  % 一个归一化模块读入所有得分。
  \begin{scope}[shift={(0,8)}]
    \node[nn transform,minimum width=22mm,minimum height=17mm] (softmax) at (114,0) {Softmax};
    \foreach \i/\dy in {1/7,2/0,3/-7}{
      \node[head unit] (c\i) at (86,\dy) {$a_{\i}$};
      \draw[nn link,draw=NNHidden!65] (79,0)--(c\i.west);
      \draw[head link] (c\i.east)--(103,\dy);
    }
    \draw[head link] (softmax.east)--(132,0);
    \node[result] at (136,3) {$\sum_k p_k=1$};
    \node[result,text=BookMuted] at (136,-3) {多类交叉熵};
  \end{scope}

  % 三个独立变换模块分别读取对应得分。
  \begin{scope}[shift={(0,-24)}]
    \foreach \i/\dy in {1/7,2/0,3/-7}{
      \node[head unit] (d\i) at (86,\dy) {$a_{\i}$};
      \draw[nn link,draw=NNHidden!65] (79,0)--(d\i.west);
      \node[nn transform,minimum width=22mm,minimum height=5mm,inner sep=1pt]
        (sig\i) at (114,\dy) {$s(a_{\i})$};
      \draw[head link] (d\i.east)--(sig\i.west);
      \draw[head link] (sig\i.east)--(132,\dy);
    }
    \node[result] at (136,3) {$p_k\in(0,1)$};
    \node[result,text=BookMuted] at (136,-3) {逐标签交叉熵};
  \end{scope}

  % 表示由下游任务提供学习约束。
  \begin{scope}[shift={(0,-56)}]
    \node[nn transform,minimum width=22mm,minimum height=17mm] (rep) at (114,0)
      {恒等或\\长度归一化};
    \foreach \i/\dy in {1/7,2/0,3/-7}{
      \node[head unit] (e\i) at (86,\dy) {$a_{\i}$};
      \draw[nn link,draw=NNHidden!65] (79,0)--(e\i.west);
      \draw[head link] (e\i.east)--(103,\dy);
    }
    \node[nn transform,draw=NNGradient,fill=NNGradient!18,
      minimum width=29mm] (downstream) at (152,0) {下游模块};
    \draw[head link] (rep.east)--(downstream.west);
    \node[text=NNOutput] at (131,7) {$\bm{r}$};
    \node[text=BookMuted,inner sep=1pt] (objective) at (152,-14) {任务损失$\ell$};
    \draw[head link,draw=NNGradient] (downstream.south)--(objective.north);
  \end{scope}
\end{tikzpicture}

  \caption{共同的前馈结构与五种可替换的输出头。左侧蓝色输入单元经紫色隐藏单元形成$\bm{h}^{(L)}$；汇合线传递整组隐藏坐标，不执行求和。右侧青色单元产生各任务自己的原始得分，每个单元都读取全部隐藏坐标，并具有独立的输出权重与偏置。(a)、(b)只需一个得分；(c)联合归一化，(d)逐坐标变换；(e)产生由下游任务约束的表示。向量头仅画三个坐标示意，实际维度由任务决定。右侧文字区分预测量与训练目标，隐藏激活无需与输出变换相同。}
  \label{fig:ffn-output-heads}
\end{figure*}

\subsubsection{回归：预测连续数值}
\label{sec:ffn-output-regression}

回归的目标是连续数值，例如经过单位缩放的房价。对标量目标$y\in\mathbb{R}$，令$q=1$，输出层只产生一个得分$a$；采用恒等变换$g(a)=a$时，预测值就是$\widehat y=a$。平方损失衡量预测值与目标的距离：
\begin{equation}
  \ell(a,y)=\frac12(a-y)^2,
  \qquad \frac{\partial\ell}{\partial a}=a-y\text{。}
  \label{eq:ffn-regression-loss}
\end{equation}
系数$1/2$用于简化导数，不改变仅有这一损失时的最优预测。若教学样本的目标为$2.5$、当前预测为$2$，损失为$0.125$，对输出的导数为$-0.5$；沿负梯度方向调整输出，会让它增大并接近目标。多个连续目标可以使用向量输出，以及$\frac12\lVert\bm{a}-\bm{y}\rVert_2^2$；不同坐标的单位与尺度需要在数据处理或损失权重中考虑。

平方损失还规定了预测的统计含义。在二阶矩有限、可选函数足够自由时，使条件期望损失最小的预测是$\mathbb{E}[y\mid\bm{x}]$，即给定输入后的条件均值。因而同一个输入对应多种可能结果时，平方损失倾向于预测它们的均值，并不直接输出所有可能结果。

输出范围应与目标配合。若预测正尺度，可以令$\widehat y=\operatorname{Softplus}(a)$，再对$\widehat y$计算损失；此时对$a$求导还需乘输出变换的导数。若同时预测不确定性，可以让网络输出均值和正方差，再用所选概率分布的负对数似然训练。约束预测值的范围与建立一个完整的条件分布，是两个不同的设计决定。

\subsubsection{二分类：预测一个正类概率}
\label{sec:ffn-output-binary}

二分类判断一个事件是否成立，例如一封邮件是否属于垃圾邮件。令$y\in\{0,1\}$，其中$1$表示正类。输出层仍只需要一个实数$a$，经过Sigmoid得到
\[
  p=s(a),\qquad P_{\bm{\theta}}(y=1\mid\bm{x})=p,
  \qquad P_{\bm{\theta}}(y=0\mid\bm{x})=1-p\text{。}
\]
这里$p$是模型给出的正类概率，$a$称为\term{logit}，即对数几率$\log(p/(1-p))$。两个类别的概率已经由$p$和$1-p$共同确定，无须再设置两个彼此独立的概率输出。

在这个伯努利模型下，正确标签的概率为$p^y(1-p)^{1-y}$。对它取负对数，就得到二元交叉熵：
\begin{equation}
  \begin{aligned}
    \ell(a,y)
      &=-y\log p-(1-y)\log(1-p)\\
      &=\log(1+\exp(a))-ya\text{。}
  \end{aligned}
  \label{eq:ffn-binary-cross-entropy}
\end{equation}
若$y=1$而$p=0.8$，损失约为$0.2231$；若把这个正例的概率降到$0.2$，损失会增至约$1.6094$。交叉熵不仅考察最后选对哪一类，还区分模型给正确标签分配了多少概率。

式\eqref{eq:ffn-binary-cross-entropy}以原始得分$a$为输入，把Sigmoid和损失合在一起计算。其稳定形式为$\max\{a,0\}-ya+\log(1+\exp(-|a|))$，导数为$\partial\ell/\partial a=p-y$。这避免了先得到极接近$0$或$1$的概率再取对数。预测类别时还需选定阈值$\tau$，例如按$p\geq\tau$判为正类；误判代价或业务约束不同，所需阈值也可能不同。硬判定用于决策，连续概率与损失用于学习。

图~\ref{fig:ffn-output-heads}(a)、(b)把回归与二分类放在同一隐藏结构上比较。两者都可以只产生一个原始得分，差别发生在输出解释和学习目标；二分类的概率变换不会要求隐藏层也使用Sigmoid。


\subsubsection{互斥多分类：预测一个类别}
\label{sec:ffn-output-multiclass}

若一个样本恰好属于$K$个类别中的一个，例如一张数字图像对应一个数字类别，需要输出一组相互竞争的概率。设$y\in\{1,\ldots,K\}$、$q=K$，输出层的每个坐标$a_k$是对应类别的原始得分。Softmax将整组得分转换为概率：
\begin{equation}
  p_k=\frac{\exp(a_k)}{\sum_{j=1}^{K}\exp(a_j)},
  \qquad \sum_{k=1}^{K}p_k=1,\qquad
  \widehat y=\operatorname*{arg\,max}_{k}p_k\text{。}
  \label{eq:ffn-multiclass-output}
\end{equation}
最大概率并列时还需约定决策规则。与逐坐标使用Sigmoid相比，Softmax的分母让所有类别共同参与归一化：提高一个类别的概率，必然改变至少一个其他类别的概率。

正确标签的概率为$p_y$。多类交叉熵就是这个概率的负对数，把Softmax代入后可写成
\begin{equation}
  \ell(\bm{a},y)=-\log p_y
  =-a_y+\log\sum_{j=1}^{K}\exp(a_j)\text{。}
  \label{eq:ffn-multiclass-cross-entropy}
\end{equation}
对$a_k$求导得到$p_k-\mathbb{I}\{y=k\}$。例如教学样本有三类，$y=2$、$\bm{p}=(0.1,0.7,0.2)^{\top}$，则损失约为$0.3567$，得分梯度为$(0.1,-0.3,0.2)^{\top}$。负梯度方向倾向于提高正确类别得分、降低其他类别得分；是否以及怎样改变这些得分，还取决于网络参数和输入。

向全部得分加同一个常数，不会改变概率。因此计算时可减去$c=\max_j a_j$，使指数不超过$1$；对数求和项计算为$c+\log\sum_j\exp(a_j-c)$。分类由相对得分决定，原始得分本身既不要求为正，也不要求加起来等于$1$。

\subsubsection{多标签分类：预测多个二值结果}
\label{sec:ffn-output-multilabel}

一个对象也可能同时具有多个标签。例如同一文档可以同时属于“技术”和“教育”，而不属于“体育”。这时要分别判断每个标签是否成立，而不能强迫模型只选择一个类别。\term{多标签分类}（multi-label classification）的目标为$\bm{y}\in\{0,1\}^{K}$，输出层产生$K$个得分，对每个坐标分别使用Sigmoid：
\begin{equation}
  \begin{aligned}
    p_k&=s(a_k),\\
    \ell(\bm{a},\bm{y})
      &=\sum_{k=1}^{K}\bigl[-y_k\log p_k-(1-y_k)\log(1-p_k)\bigr]\\
      &=\sum_{k=1}^{K}\bigl[\log(1+\exp(a_k))-y_k a_k\bigr]\text{。}
  \end{aligned}
  \label{eq:ffn-multilabel-loss}
\end{equation}
每个坐标都对应一次二分类，概率总和不必为$1$。例如目标为$(1,1,0)^{\top}$，预测为$(0.8,0.7,0.1)^{\top}$，概率总和为$1.6$，按阈值$0.5$会选出前两个标签。损失中前两项奖励相应概率增大，第三项则奖励该标签的概率减小。

式\eqref{eq:ffn-multilabel-loss}采用标签损失求和；也可以按标签数取平均，但会相应改变梯度尺度。每个得分的导数仍为$p_k-y_k$，不同标签可以设置各自的决策阈值。这里“分别使用”描述输出变换和损失的组织方式；若把各标签的伯努利概率乘起来解释为联合分布，就采用了给定输入后的标签条件独立模型。共享隐藏层能利用共同特征，但这种输出头本身没有显式建立标签之间的联合依赖。

图~\ref{fig:ffn-output-heads}(c)、(d)对比了同样输出三个得分的两种任务。Softmax把三个类别放进一次竞争，逐坐标Sigmoid则允许多个标签同时成立。输出维度相同，不能据此选择相同的概率变换和决策规则。


\subsubsection{表示输出：由下游目标提供约束}
\label{sec:ffn-output-representation}

有时网络的输出不会立即解释成类别或数值，而是交给另一个模块做比较、分类或重建。例如，编码器输出一组坐标，后续分类头再根据这些坐标预测标签。此时可以直接使用$\bm{r}=\bm{a}\in\mathbb{R}^{q}$，也可以根据下游任务约束表示的长度；当$\bm{a}\neq\bm{0}$时，单位长度表示为
\begin{equation}
  \bm{r}=\frac{\bm{a}}{\lVert\bm{a}\rVert_2}\text{。}
  \label{eq:ffn-normalized-representation}
\end{equation}
例如$(3,4)^{\top}$会变成$(0.6,0.8)^{\top}$。这一变换保留方向，去掉整体长度差异；它改变了输出的几何约束，却没有指定哪些对象应当相近、哪些坐标应当携带什么信息。

表示的学习信号必须来自目标。接上分类头时，分类损失的梯度经过这个头传回表示网络；做重建时，由重建误差规定需要保留的信息；按样本关系学习时，则需说明怎样奖励相近样本、区分不同样本。只要求相近样本的表示相同，可能让全部输入得到同一个表示；单位长度约束也不能排除这种情况。因而“输出一个向量”与“学到有用表示”之间，还缺少与使用目的相配的训练约束。

图~\ref{fig:ffn-output-heads}(e)把表示输出与下游目标相连。前馈网络产生表示，下游模块把表示用于具体任务，损失的梯度再沿这条计算路径传回。若去掉下游目标，就只剩一次向量变换，还没有规定应学到什么。


\subsubsection{从单样本损失到训练目标}
\label{sec:ffn-empirical-objective}

上面的输出头规定了一条预测怎样计分。训练时，所有样本共享同一组参数，目标还需汇总各样本损失。设训练集为$S=((\bm{x}_i,\bm{y}_i))_{i=1}^{n}$，标量标签按对应任务书写，记网络对第$i$个输入产生的原始输出为$\bm{a}_i(\bm{\theta})$，则
\begin{equation}
  J_S(\bm{\theta})=
  \frac1n\sum_{i=1}^{n}\ell\!\left(\bm{a}_i(\bm{\theta}),\bm{y}_i\right)
  +\lambda\Omega(\bm{\theta})\text{。}
  \label{eq:ffn-training-objective}
\end{equation}
这里把输出变换并入损失，回归中的恒等变换、二分类的Sigmoid和多分类的Softmax分别按上面的公式处理。表示输出则把下游目标与表示网络连成同一计算路径。$\Omega$是可选的参数惩罚，$\lambda\geq0$控制惩罚强度。

输出维度、隐藏层数和层宽是设计时确定或用验证集选择的超参数，权重与偏置则从训练数据学习。概率输出满足取值范围，只说明模型形式正确；预测概率是否与实际发生频率一致，还需用数据评价。由此，输出头回答“网络预测什么”，损失回答“训练奖励什么”，决策规则回答“预测怎样用于行动”。这三者应当在同一任务设定下配合。

\section{参数学习}
\label{sec:ffn-learning}

确定结构和目标之后，学习过程需要一个起点，也需要知道每个参数应沿哪个方向调整。初始化影响初始表示及梯度的尺度，反向传播计算各参数对损失的局部影响，自动微分把这些计算规则应用到程序上。优化方法使用所得梯度更新参数；反向传播本身并不负责选择学习率或保证找到全局最优解。

\subsection{参数初始化}
\label{sec:ffn-initialization}

初始化要解决两个问题：让同层单元从不同起点学习，并让数值与梯度在层间传播时保持合适尺度。Sigmoid、Tanh和ReLU都提供非线性，但对输入的压缩方式不同，因此随机权重不能只按层宽设置，还需与隐藏激活配合。下面先给出共同的尺度估计，再按饱和响应、整流响应和其他用途分别讨论。

\subsubsection{共同要求：打破对称并控制尺度}
\label{sec:ffn-initialization-common}

若同一隐藏层单元的输入权重、偏置以及相应的输出连接完全相同，在相同数据和对称更新规则下，它们的输出与梯度也会保持相同。全零初始化就是典型的对称起点：增加了单元数，却没有让它们承担不同的计算。随机初始化权重可以打破这种对称性；偏置通常可初始化为零，因为不同权重已经让单元响应有所区别。

随机性的作用不等于尺度控制。设某层输入坐标的二阶矩为$q_{\mathrm{in}}=\mathbb{E}[u_k^2]$，初始权重独立、均值为零、方差为$v$，权重与输入独立，偏置为零。对有$m$个输入的单元，$z=\sum_{k=1}^{m}w_k u_k$满足
\begin{equation}
  \mathbb{E}[z^2]=m v q_{\mathrm{in}}\text{。}
  \label{eq:ffn-initial-signal-scale}
\end{equation}
交叉项因不同权重独立且均值为零而消失。若激活保持输入尺度，则$v$应随$m$缩小，避免连接数增加时预激活幅度持续增大；若激活会压缩或截断输入，还需补上相应影响。

输入连接数$m=d_{\ell-1}$称为\term{fan-in}；同一输入坐标通向的输出连接数$d_\ell$称为\term{fan-out}。前向计算关心每个输出汇总了多少输入，后向计算关心每个输入汇总了多少输出的影响。两种连接数不同，就可能需要在保持前向尺度与保持后向尺度之间作折中。

\subsubsection{Sigmoid与Tanh：同一饱和族，中心尺度不同}
\label{sec:ffn-initialization-saturating}

本章的Sigmoid特指Logistic函数$s(z)$。它与Tanh有精确关系：
\begin{equation}
  \tanh(z)=2s(2z)-1,\qquad
  s(z)=\frac{1+\tanh(z/2)}{2}\text{。}
  \label{eq:ffn-sigmoid-tanh-relation}
\end{equation}
因此两者属于相近的S形饱和响应：足够大的正负输入都会进入变化缓慢的尾部。它们的区别在于坐标缩放与输出中心。Sigmoid的输出位于$(0,1)$，$s(0)=1/2$；Tanh的输出位于$(-1,1)$，$\tanh(0)=0$。在关于零对称的输入分布下，Tanh因奇对称性具有零均值，而Sigmoid的均值为$1/2$；非对称输入下则不能保证Tanh输出零均值。

两者在中心附近的斜率也不同：$s'(0)=1/4$，$\tanh'(0)=1$。如果许多隐藏层在中心附近使用Sigmoid，激活导数本身就会反复缩小反向信号；若预激活幅度过大，Sigmoid和Tanh又都会饱和。增大权重虽然能放大仿射变换，却可能把更多单元推入小导数区域，不能靠无限放大权重补偿衰减。

对中心附近近似单位斜率的激活，把式\eqref{eq:ffn-initial-signal-scale}用于前向传播，得到$v\approx1/\mathrm{fan\text{-}in}$；对后向传播作类似独立性近似，得到$v\approx1/\mathrm{fan\text{-}out}$。\term{Glorot初始化}（Glorot initialization，也称Xavier初始化）采用折中方差
\begin{equation}
  \begin{aligned}
    v&=\frac{2}{d_{\ell-1}+d_\ell},\\
    (\bm{W}_\ell)_{j,k}&\sim
    \mathcal{U}\!\left[-\sqrt{\frac{6}{d_{\ell-1}+d_\ell}},
    \sqrt{\frac{6}{d_{\ell-1}+d_\ell}}\right]\text{。}
  \end{aligned}
  \label{eq:ffn-glorot-initialization}
\end{equation}
$\mathcal{U}$表示均匀分布，也可以使用同方差的零均值正态分布。Tanh在原点附近符合上述单位斜率近似，因而常与这一尺度配合；实际设置可再按所选激活的增益调整。Glorot与Bengio的原文同时分析了Sigmoid的非零输出中心和饱和问题，这些因素会使近似线性传播的条件更难成立\scite{glorot2010initialization}

\subsubsection{ReLU及其变体：补偿被截断的响应}
\label{sec:ffn-initialization-rectifiers}

ReLU采用不同的响应机制：它在正半轴保留输入、导数为$1$，在负半轴输出和导数都为$0$。因此它没有正端饱和，但会截断一部分信号；若单元对全部训练样本都处于负区间，又可能长期没有输入参数的学习信号。前述饱和函数两端导数趋零，与这里负区间导数严格为零，是不同的机制。

若预激活$z$的分布关于零对称，ReLU满足
\[
  \mathbb{E}[\rho(z)^2]=\frac12\mathbb{E}[z^2]\text{。}
\]
这条关系控制的是二阶矩；ReLU输出通常具有正均值，不能直接称为方差保持。把截断因子$1/2$代入式\eqref{eq:ffn-initial-signal-scale}，为近似保持前向激活的二阶矩，\term{He初始化}（He initialization，也称Kaiming初始化）采用
\begin{equation}
  (\bm{W}_\ell)_{j,k}\sim
  \mathcal{N}\!\left(0,\frac{2}{d_{\ell-1}}\right)\text{。}
  \label{eq:ffn-he-initialization}
\end{equation}
相对于单位增益响应的$1/\mathrm{fan\text{-}in}$，系数$2$补偿了对称分布下负半轴被截断的二阶矩。

若负半轴保留斜率$\alpha$，则$\mathbb{E}[\sigma(z)^2]=(1+\alpha^2)\mathbb{E}[z^2]/2$，对应前向权重方差为$2/((1+\alpha^2)d_{\ell-1})$。这适用于初始化时固定斜率的Leaky ReLU，也可用于PReLU的初始斜率；斜率开始学习后，初始尺度不再是全过程保证。若主要保持后向梯度的二阶矩，可将分母中的fan-in换成fan-out，所控制的传播方向应明确\scite{he2015rectifiers}

\subsubsection{其他隐藏激活与任务输出：分别设置}
\label{sec:ffn-initialization-other}

GELU、SiLU和ELU的负区间响应不同，不具有ReLU完全相同的截断关系，因此不能仅凭正端近似线性就原样套用系数$2$。SELU又是把激活常数与特定统计条件配套考虑的设计，应沿用其对应的初始化和结构条件。选择其他激活时，需要检查所用实现的增益约定、输入尺度和层宽，并观察初始活动值与梯度；这些检查没有给出统一的训练全过程保证。

隐藏激活与输出变换也应分开判断。二分类网络可以在隐藏层用ReLU、在输出处用Sigmoid：前者构造表示，后者把一个得分解释成概率，它们承担不同职责。ReLU隐藏层可采用式\eqref{eq:ffn-he-initialization}，但线性得分层没有负半轴截断，无须自动继承同一系数。二分类输出还需关注初始logit是否过大，而不能把“最后使用Sigmoid”理解为所有隐藏层都要按饱和函数初始化。

以上尺度推导利用的是初始化时的分布与独立性假设。训练后，权重、表示与上游梯度可能相关，因而保持初始二阶矩不等于一直保持梯度，也不等于保证优化成功。初始化、输入标准化、激活函数与更新方式需要共同配合；后面的反向传播将说明梯度具体如何穿过这些计算。

\subsection{反向传播}
\label{sec:ffn-backpropagation}

隐藏层没有直接给定的目标值，怎样判断某个隐藏权重应当增大还是减小？一个权重改变了本层输出，本层输出改变了后续表示，后续表示又改变最终损失。链式法则把这些局部影响相乘并汇总。\term{反向传播}（backpropagation）按照与前向依赖相反的顺序组织这项计算，让不同参数重复利用已经算出的中间导数。

先固定一个样本，把输出变换与损失合为标量函数$\ell(\bm{a},\bm{y})$。用列向量表示梯度，定义
\begin{equation}
  \bm{\delta}^{(L+1)}=\nabla_{\bm{a}}\ell,
  \qquad
  \bm{\delta}^{(\ell)}=\nabla_{\bm{z}^{(\ell)}}\ell
  \quad(\ell=1,\ldots,L)\text{。}
  \label{eq:ffn-delta-definition}
\end{equation}
向量$\bm{\delta}^{(\ell)}\in\mathbb{R}^{d_\ell}$记录损失对本层各预激活坐标的敏感程度。它常被称为反向传播的“误差”，但一般并不是本层输出减去某个隐藏标签；隐藏层并没有这样的标签。

由式\eqref{eq:ffn-output-map}，输出得分对$\bm{h}^{(L)}$的雅可比为$\bm{W}_{L+1}$，所以损失对最后一个隐藏表示的梯度为$\bm{W}_{L+1}^{\top}\bm{\delta}^{(L+1)}$。隐藏激活逐坐标作用，其雅可比是以各坐标导数为对角元素的矩阵。把两步链式法则合起来，便得到逐层递推：
\begin{equation}
  \bm{\delta}^{(\ell)}=
  \left(\bm{W}_{\ell+1}^{\top}\bm{\delta}^{(\ell+1)}\right)
  \odot\sigma_\ell'\!\left(\bm{z}^{(\ell)}\right),
  \qquad \ell=L,\ldots,1\text{。}
  \label{eq:ffn-backward-recursion}
\end{equation}
其中$\odot$表示逐坐标乘法。矩阵转置把后层各单元的影响汇总到当前坐标，激活导数再调整这些影响能有多少通过本层。递推中的每个向量维度都是$d_\ell$，因此可逐层核对计算是否匹配。

参数的导数来自本层仿射运算。对任一层$\ell=1,\ldots,L+1$，第$j$个预激活值对权重$(\bm{W}_\ell)_{j,k}$的导数是$h_k^{(\ell-1)}$，对偏置$(\bm{b}_\ell)_j$的导数是$1$。由此得到
\begin{equation}
  \nabla_{\bm{W}_\ell}\ell
    =\bm{\delta}^{(\ell)}(\bm{h}^{(\ell-1)})^{\top},
  \qquad
  \nabla_{\bm{b}_\ell}\ell=\bm{\delta}^{(\ell)}\text{。}
  \label{eq:ffn-parameter-gradients}
\end{equation}
第一式是一个外积：每条连接的梯度等于目标单元的损失敏感度乘以来源单元的活动值。矩阵形状为$d_\ell\times d_{\ell-1}$，恰好与权重相同。对输入的梯度则为$\nabla_{\bm{x}}\ell=\bm{W}_1^{\top}\bm{\delta}^{(1)}$；即使输入不参与参数更新，这个量也能用于分析预测对输入变化的敏感程度。

输出层的起点由任务决定。采用式\eqref{eq:ffn-regression-loss}时，$\delta^{(L+1)}=a-y$；采用二元交叉熵时，$\delta^{(L+1)}=p-y$；采用多类交叉熵时，第$k$个坐标为$p_k-\mathbb{I}\{y=k\}$，其中$\mathbb{I}$是事件成立时为$1$的指示函数。概率输出与交叉熵合并后得到的简洁导数，是两者配合的结果。若Sigmoid输出后使用平方损失，导数还包含$s'(a)$，不能沿用$p-y$。

图~\ref{fig:ffn-backpropagation}说明了前向值与反向梯度的区别。前向计算提供各层的预激活和表示；反向计算从损失开始，重复利用这些值与局部导数。梯度向上游传播不会改变本次前向预测已经使用的输入，也不会给预测网络增加有向环。

\begin{figure*}[htb]
  \centering
  % 原创教学示意；z 与 h 是同一隐藏层激活前后的坐标，并非两个隐藏层。
\begin{tikzpicture}[x=1mm,y=1mm,font=\figurefont\footnotesize,text=NNInk,
  bp neuron/.style={circle,minimum size=7mm,inner sep=0pt,line width=1.1pt},
  bp input/.style={bp neuron,draw=NNInput,fill=NNInput!18},
  bp hidden/.style={bp neuron,draw=NNHidden,fill=NNHidden!18},
  bp output/.style={bp neuron,draw=NNOutput,fill=NNOutput!18},
  bp link/.style={draw=NNLine,line width=.85pt},
  bp forward/.style={-{Stealth[length=1.7mm,width=1.2mm]},draw=NNHidden,
    line width=.85pt},
  bp backward/.style={-{Stealth[length=1.8mm,width=1.3mm]},draw=NNGradient,
    line width=1.2pt,dashed},
  bp adjoint/.style={text=NNGradient,inner sep=2pt}]
  \foreach \x/\title/\symbol in {
    10/输入/\bm{x},47/预激活/\bm{z}^{(1)},84/隐藏表示/\bm{h}^{(1)},
    121/输出得分/\bm{a},156/损失/\ell}{
    \node[font=\figurefont\bfseries] at (\x,29) {\title};
    \node at (\x,21) {$\symbol$};
  }
  % 先画连接，再画节点，使细线在神经元边缘干净地结束。
  \foreach \iy in {9,0,-9}
    \foreach \hy in {13.5,4.5,-4.5,-13.5}
      \draw[bp link] (10,\iy)--(47,\hy);
  \foreach \hy in {13.5,4.5,-4.5,-13.5}
    \foreach \oy in {6,-6}
      \draw[bp link] (84,\hy)--(121,\oy);
  \foreach \i/\y in {1/9,2/0,3/-9}
    \node[bp input] (x\i) at (10,\y) {$x_{\i}$};
  \foreach \i/\y in {1/13.5,2/4.5,3/-4.5,4/-13.5}{
    \node[bp hidden,fill=white] (z\i) at (47,\y) {$z_{\i}$};
    \node[bp hidden] (h\i) at (84,\y) {$h_{\i}$};
    \draw[bp forward] (z\i)--(h\i);
  }
  \foreach \i/\y in {1/6,2/-6}
    \node[bp output] (a\i) at (121,\y) {$a_{\i}$};
  \node[circle,draw=NNGradient,fill=NNGradient!18,line width=1.1pt,
    minimum size=10mm,inner sep=0pt] (loss) at (156,0) {$\ell$};
  \draw[bp forward,draw=NNOutput] (a1)--(loss.north west);
  \draw[bp forward,draw=NNOutput] (a2)--(loss.south west);
  \node[text=NNInput] at (28.5,15) {$\bm{W}_1,\bm{b}_1$};
  \node[text=NNHidden] at (65.5,15) {$\sigma_1$};
  \node[text=NNHidden] at (102.5,15) {$\bm{W}_2,\bm{b}_2$};
  \node[text=NNOutput] at (140,15) {$\ell(\bm{a},y)$};
  \draw[NNLine,line width=.45pt] (1,-20)--(164,-20);
  \node[bp adjoint] (b1) at (10,-35) {$\nabla_{\bm{x}}\ell$};
  \node[bp adjoint] (b2) at (47,-35) {$\bm{\delta}^{(1)}$};
  \node[bp adjoint] (b3) at (84,-35) {$\nabla_{\bm{h}^{(1)}}\ell$};
  \node[bp adjoint] (b4) at (121,-35) {$\bm{\delta}^{(2)}$};
  \node[bp adjoint] (b5) at (156,-35) {$1$};
  \draw[bp backward] (b2)--node[above] {$\bm{W}_1^{\top}$} (b1);
  \draw[bp backward] (b3)--node[above] {$\odot\sigma_1'$} (b2);
  \draw[bp backward] (b4)--node[above] {$\bm{W}_2^{\top}$} (b3);
  \draw[bp backward] (b5)--node[above] {$\nabla_{\bm{a}}\ell$} (b4);
  \draw[bp backward] (b2.south)--(47,-47);
  \draw[bp backward] (b4.south)--(121,-47);
  \node[align=center,inner sep=0pt] at (47,-55)
    {$\nabla_{\bm{W}_1}\ell=\bm{\delta}^{(1)}\bm{x}^{\top}$\\[2pt]
     $\nabla_{\bm{b}_1}\ell=\bm{\delta}^{(1)}$};
  \node[align=center,inner sep=0pt] at (121,-55)
    {$\nabla_{\bm{W}_2}\ell=\bm{\delta}^{(2)}(\bm{h}^{(1)})^{\top}$\\[2pt]
     $\nabla_{\bm{b}_2}\ell=\bm{\delta}^{(2)}$};
  \draw[bp forward,draw=NNInput] (28,-70)--(38,-70);
  \node[anchor=west] at (41,-70) {前向计算};
  \draw[bp backward] (99,-70)--(89,-70);
  \node[anchor=west] at (102,-70) {梯度回传};
\end{tikzpicture}

  \caption{一个隐藏层上的前向值与反向梯度。圆点展示向量坐标，预激活与隐藏表示两列对应同一隐藏层，紫色箭头表示逐坐标激活；细线表示仿射层的全连接。下方赭色虚线从损失对自身的导数$1$开始回传，与上方变量逐列对应，并标出局部求导规则。$\sigma_1'$在保存的$\bm{z}^{(1)}$处计算。底部由敏感度与前层活动值形成参数梯度，待所有梯度计算完成后才更新参数。}
  \label{fig:ffn-backpropagation}
\end{figure*}

\begin{example}[一次前向与反向计算]
  \label{ex:ffn-numerical-backpropagation}
  沿用异或的二值输入，取一个隐藏层、两个ReLU单元及一个二分类输出。下面的参数是教学构造：
  \[
    \begin{aligned}
      \bm{W}_1&=\begin{pmatrix}1&-1\\-1&1\end{pmatrix},
      &\bm{b}_1&=\begin{pmatrix}0.2\\0.2\end{pmatrix},\\
      \bm{W}_2&=\begin{pmatrix}1&1\end{pmatrix},
      &b_2&=-0.5\text{。}
    \end{aligned}
  \]
  对样本$\bm{x}=(1,0)^{\top}$、$y=1$，前向结果为
  \[
    \bm{z}^{(1)}=(1.2,-0.8)^{\top},\quad
    \bm{h}^{(1)}=(1.2,0)^{\top},\quad a=0.7\text{。}
  \]
  因而$p=s(0.7)\approx0.6682$，二元交叉熵$\ell=-\log p\approx0.4032$。非零偏置让本例的两个预激活值都远离ReLU的不可微点。

  输出敏感度为$\delta^{(2)}=p-1\approx-0.3318$。第一个隐藏单元处于正半轴，第二个处于负半轴，于是
  \[
    \begin{aligned}
      \bm{\delta}^{(1)}&\approx(-0.3318,0)^{\top},\\
      \nabla_{\bm{W}_1}\ell&\approx
      \begin{pmatrix}-0.3318&0\\0&0\end{pmatrix},
      &\nabla_{\bm{b}_1}\ell&\approx
      \begin{pmatrix}-0.3318\\0\end{pmatrix}\text{。}
    \end{aligned}
  \]
  输出参数的梯度为$\nabla_{\bm{W}_2}\ell\approx(-0.3982,0)$、$\partial\ell/\partial b_2\approx-0.3318$。第一层第一个单元对$x_2$的权重梯度为零，是因为这个样本的$x_2=0$；第二个单元的梯度为零，则来自负半轴上的激活导数。两种零梯度有不同原因。

  若仅用该样本、不加正则项，并以学习率$0.1$同时更新所有参数，新的原始得分约为$0.8500$，正类概率约为$0.7006$，损失约为$0.3559$。这一次更新提高了该正例的概率；它不能说明四个异或样本的总损失也一定下降，更不能说明新数据上的风险已经降低。
\end{example}

训练目标通常包含多个样本，它们共享同一组参数，因此梯度需要相加。若本次更新读取一个索引集合$B\subseteq\{1,\ldots,n\}$，称为\term{小批量}（mini-batch），其平均数据梯度为
\begin{equation}
  \bm{g}_\ell=
  \frac1{|B|}\sum_{i\in B}
  \bm{\delta}_i^{(\ell)}(\bm{h}_i^{(\ell-1)})^{\top},
  \qquad
  \bm{r}_\ell=
  \frac1{|B|}\sum_{i\in B}\bm{\delta}_i^{(\ell)}\text{。}
  \label{eq:ffn-minibatch-gradients}
\end{equation}
这里$\bm{g}_\ell$和$\bm{r}_\ell$分别对应权重与偏置，每个样本的$\bm{\delta}_i^{(\ell)}$按未除以批量大小的单样本损失计算。若反向起点已经使用批平均损失，便不应在这里再次除以$|B|$。矩阵实现可将样本按列排列，用同样的矩阵乘法一次计算这些外积之和；样本按行排列的实现会相应改变转置位置。

若$\Omega=\frac12\sum_\ell\lVert\bm{W}_\ell\rVert_{\mathrm{F}}^2$，其中$\lVert\cdot\rVert_{\mathrm{F}}$是矩阵各元素平方和开方得到的Frobenius范数，则权重总梯度为$\bm{g}_\ell+\lambda\bm{W}_\ell$，未被惩罚的偏置不加这一项。给定学习率$\eta_t>0$，普通梯度更新为
\begin{equation}
  \begin{aligned}
    \bm{W}_\ell^{(t+1)}
      &=\bm{W}_\ell^{(t)}-\eta_t(\bm{g}_\ell+\lambda\bm{W}_\ell^{(t)}),\\
    \bm{b}_\ell^{(t+1)}&=\bm{b}_\ell^{(t)}-\eta_t\bm{r}_\ell\text{。}
  \end{aligned}
  \label{eq:ffn-gradient-update}
\end{equation}
这里迭代上标$t$与层下标$\ell$分别记录训练时刻和结构位置。

\begin{algorithm}[前馈网络的小批量梯度学习]
  \label{alg:ffn-minibatch-learning}
  \begin{algorithmic}[1]
    \Require 训练集$S$、验证集、网络结构、损失、学习率序列与停止规则
    \Ensure 经验证集选择的参数$\bm{\theta}$
    \State 按激活与层宽初始化权重，初始化偏置
    \While{尚未达到停止规则}
      \State 读取小批量$B$，保存各样本的预激活、隐藏表示与输出
      \State 计算输出敏感度$\bm{\delta}_i^{(L+1)}$
      \For{$\ell=L,L-1,\ldots,1$}
        \State 按式\eqref{eq:ffn-backward-recursion}计算$\bm{\delta}_i^{(\ell)}$
      \EndFor
      \State 汇总所有层的数据梯度，并加入所选惩罚项的梯度
      \State 使用同一更新前参数计算的梯度，同时更新各层参数
      \State 按既定间隔评估验证集，记录待选择的参数
    \EndWhile
    \State 返回按预先规定验证指标选择的参数
  \end{algorithmic}
\end{algorithm}

算法~\ref{alg:ffn-minibatch-learning}把求导和更新分成两步。如果在尚未完成反向计算时就改写某层权重，前后各层的梯度可能对应不同的参数，便不再是同一个目标在同一点的梯度。停止规则可以包含最大更新次数和验证指标的停止条件；优化器、学习率安排及更广泛的训练问题将在训练章比较。

反向传播的计算量与前向矩阵运算同阶，因为每层主要计算转置矩阵乘法和外积，而不是为每个参数重新运行网络。它通常还要保存各层的活动值，单样本隐藏激活的存储量与$\sum_\ell d_\ell$同阶，小批量时随批量大小增长。逐层导数相乘也解释了小梯度和大梯度如何累积：能够正确求出梯度，并不保证这个梯度具有适合更新的尺度。

\subsection{自动梯度计算}
\label{sec:ffn-autodiff}

普通串联网络的梯度可以手工推导，但模型加入多个分支、共享参数或额外损失以后，逐项维护公式很容易遗漏一条影响路径。要把求导交给程序，需要说明程序到底在算什么。数值微分、符号微分和自动微分都能提供导数信息，采用的原理却不同。

\term{数值微分}（numerical differentiation）通过改变输入后观察函数值之差来近似导数。例如，用第$j$个标准基向量$\bm{e}_j$和步长$\varepsilon>0$，中心差分为
\begin{equation}
  \frac{\partial J}{\partial\theta_j}
  \approx
  \frac{J(\bm{\theta}+\varepsilon\bm{e}_j)
       -J(\bm{\theta}-\varepsilon\bm{e}_j)}{2\varepsilon}\text{。}
  \label{eq:ffn-finite-difference}
\end{equation}
在足够光滑且附近高阶导数受控时，其截断误差随$\varepsilon^2$缩小；步长过小时，两次接近的函数值相减又会放大浮点误差。得到$p$个参数的完整梯度通常需要约$2p$次目标计算，所以它适合检查少量梯度或方向导数，不适合直接承担大网络的每次更新。检查ReLU网络时应避开激活切换点，否则差分可能跨越两个不同的线性区域。

\term{符号微分}（symbolic differentiation）按照求导规则变换表达式，例如把$\frac12(uv+u)^2$变成关于$u,v$的导数公式，再在某一点求值。它能给出可以继续分析的表达式，但未经简化和公共子表达式复用时，复杂复合函数的导数表达式可能迅速膨胀。符号方法也可以利用共享结构；因此不能把表达式膨胀说成所有符号求导实现都无法避免的性质。

\term{自动微分}（automatic differentiation, AD）把程序实际执行的计算拆成具有已知局部导数的基本运算，在计算数值的同时或之后应用链式法则。它不通过有限扰动估计导数，也不必先展开整个导数表达式。若所执行运算可微且局部求导规则正确，所得结果是这条计算路径的导数在当前点的数值，仍受浮点舍入影响；自动微分并不意味着无限精度，也不意味着不可微的程序自动变得可微。

\subsubsection{算子图与前向模式}
\label{sec:ffn-autodiff-forward}

为了区分两种传播方向，考虑教学函数$J(u,v)=\frac12(uv+u)^2$。把计算拆成$r=uv$、$s=r+u$、$t=s^2$和$J=t/2$，就得到乘法、加法、平方与缩放四个算子。每个算子只读取前面已经得到的值，输入$u$则分别进入乘法和加法，形成两条影响输出的路径。

\term{前向模式}（forward mode）随前向值一起传播某个输入方向上的变化率。若点上的输入方向为$(\dot u,\dot v)$，有
\begin{equation}
  \begin{aligned}
    \dot r&=v\dot u+u\dot v,&\dot s&=\dot r+\dot u,\\
    \dot t&=2s\dot s,&\dot J&=\tfrac12\dot t\text{。}
  \end{aligned}
  \label{eq:ffn-forward-mode}
\end{equation}
圆点表示沿所选方向的导数，不表示时间演化。取$u=2,v=3$，前向值依次为$r=6,s=8,t=64,J=32$。如果选择$(\dot u,\dot v)=(1,0)$，表示只沿$u$的正方向考察变化：乘法给出$\dot r=3$，加法还收到$u$直接传来的$1$，于是$\dot s=4$；经过平方与缩放，得到$\dot t=64$、$\dot J=32$。

\begin{figure*}[htb]
  \centering
  % 原创算子图；两种 AD 模式共用节点位置，便于直接比较传播方向。
\begin{tikzpicture}[x=1mm,y=1mm,font=\figurefont\footnotesize,text=NNInk,
  ad input/.style={circle,draw=NNInput,fill=NNInput!18,line width=1.1pt,
    minimum size=11mm,inner sep=0pt},
  ad operator/.style={circle,draw=NNHidden,fill=NNHidden!18,line width=1.1pt,
    minimum size=11mm,inner sep=0pt},
  ad forward/.style={-{Stealth[length=1.8mm,width=1.3mm]},draw=NNInput,
    line width=1.2pt,rounded corners=2mm},
  ad derivative/.style={text=NNGradient,align=center,text width=30mm,
    inner sep=0pt}]
  \node[ad input] (u) at (12,14) {$u=2$};
  \node[ad input] (v) at (12,-14) {$v=3$};
  \node[ad derivative] at (12,3) {$\dot u=1$};
  \node[ad derivative] at (12,-25) {$\dot v=0$};
  \node[ad operator] (mul) at (52,0) {$\times$};
  \node[ad operator] (add) at (84,0) {$+$};
  \node[ad operator] (square) at (116,0) {$(\cdot)^2$};
  \node[ad operator,draw=NNOutput,fill=NNOutput!18] (half) at (148,0)
    {$\times\frac12$};
  \draw[ad forward] (u.east)--(mul.north west);
  \draw[ad forward] (v.east)--(mul.south west);
  \draw[ad forward] (u.north)--(12,28)--(84,28)--(add.north);
  \node[above,text=NNInput] at (49,28) {$u=2,\quad\dot u=1$};
  \draw[ad forward] (mul)--(add);
  \draw[ad forward] (add)--(square);
  \draw[ad forward] (square)--(half);
  \foreach \x/\value in {52/r=6,84/s=8,116/t=64,148/J=32}{
    \node[text=NNInput] at (\x,-12) {$\value$};
    \draw[NNLine,line width=.45pt] (\x-12,-20)--(\x+12,-20);
  }
  \node[ad derivative] at (52,-30) {$\begin{aligned}\dot r&=3\cdot1+2\cdot0\\&=3\end{aligned}$};
  \node[ad derivative] at (84,-30) {$\begin{aligned}\dot s&=3+1\\&=4\end{aligned}$};
  \node[ad derivative] at (116,-30) {$\begin{aligned}\dot t&=2\cdot8\cdot4\\&=64\end{aligned}$};
  \node[ad derivative] at (148,-30) {$\begin{aligned}\dot J&=\tfrac12\cdot64\\&=32\end{aligned}$};
  \draw[ad forward] (22,-47)--(32,-47);
  \node[anchor=west] at (35,-47) {计算依赖};
  \node[anchor=west,text=NNInput] at (82,-47) {$r=6$};
  \node[anchor=west] at (97,-47) {值};
  \node[anchor=west,text=NNGradient] at (111,-47) {$\dot r=3$};
  \node[anchor=west] at (127,-47) {方向导数};
\end{tikzpicture}

  \caption{具体算子图上的前向模式。圆形节点表示乘法、加法、平方和缩放，蓝色实线沿原计算依赖传播；节点下方依次列出蓝色计算值和赭色方向导数。输入方向$(1,0)$给出$\partial J/\partial u=32$；加法处的$\dot s=3+1$包含了$u$的两条路径。}
  \label{fig:ffn-autodiff-forward}
\end{figure*}

图~\ref{fig:ffn-autodiff-forward}把一次方向导数计算附在原有算子图上。若要计算对$v$的偏导，改取$(\dot u,\dot v)=(0,1)$：此时$\dot r=2$、$\dot s=2$、$\dot t=32$，最终$\dot J=16$。两次遍历的前向值相同，变化的是输入方向以及随之传播的导数。

\FloatBarrier
\subsubsection{后向模式与梯度累加}
\label{sec:ffn-autodiff-reverse}

若希望从一个标量目标同时得到多个输入的导数，可以反过来问：输出对当前中间值有多敏感？\term{后向模式}（reverse mode）从输出向输入累积这种敏感度。记$\bar r=\partial J/\partial r$，对其他变量采用相同横线记号；前向值仍保留在原算子图中，从$\bar J=1$开始回传：
\begin{equation}
  \begin{aligned}
    \bar t&=\tfrac12\bar J,&\bar s&=2s\bar t,\\
    \bar r&=\bar s,&\bar u&=\bar r v+\bar s,\\
    &&\bar v&=\bar r u\text{。}
  \end{aligned}
  \label{eq:ffn-reverse-mode}
\end{equation}
每条回传贡献等于下游敏感度乘该条依赖的局部导数。缩放算子给出$\bar t=1/2$，平方算子利用保存的$s=8$得到$\bar s=8$，加法向两个输入各传递$8$。其中传给$r$的$8$再经过乘法，分别给$u$和$v$贡献$8v=24$与$8u=16$。

\begin{figure*}[htb]
  \centering
  % 原创算子图；几何位置与前向模式相同，虚线反向传递敏感度。
\begin{tikzpicture}[x=1mm,y=1mm,font=\figurefont\footnotesize,text=NNInk,
  ad input/.style={circle,draw=NNInput,fill=NNInput!18,line width=1.1pt,
    minimum size=11mm,inner sep=0pt},
  ad operator/.style={circle,draw=NNHidden,fill=NNHidden!18,line width=1.1pt,
    minimum size=11mm,inner sep=0pt},
  ad reverse/.style={-{Stealth[length=1.8mm,width=1.3mm]},draw=NNGradient,
    line width=1.2pt,dashed,rounded corners=2mm},
  ad derivative/.style={text=NNGradient,align=center,text width=30mm,
    inner sep=0pt}]
  \node[ad input] (u) at (12,14) {$u=2$};
  \node[ad input] (v) at (12,-14) {$v=3$};
  \node[ad derivative] at (12,3) {$\bar u=32$};
  \node[ad derivative] at (12,-25) {$\bar v=16$};
  \node[ad operator] (mul) at (52,0) {$\times$};
  \node[ad operator] (add) at (84,0) {$+$};
  \node[ad operator] (square) at (116,0) {$(\cdot)^2$};
  \node[ad operator,draw=NNOutput,fill=NNOutput!18] (half) at (148,0)
    {$\times\frac12$};
  \draw[ad reverse] (mul.north west)--
    node[above,sloped,text=NNGradient] {$8\cdot3=24$} (u.east);
  \draw[ad reverse] (mul.south west)--
    node[below,sloped,text=NNGradient] {$8\cdot2=16$} (v.east);
  \draw[ad reverse] (add.north)--(84,28)--(12,28)--(u.north);
  \node[above,text=NNGradient] at (49,28) {$8\cdot1=8$};
  \draw[ad reverse] (add)--node[above,text=NNGradient] {$\times1$} (mul);
  \draw[ad reverse] (square)--node[above,text=NNGradient] {$\times2s$} (add);
  \draw[ad reverse] (half)--node[above,text=NNGradient] {$\times\tfrac12$} (square);
  \foreach \x/\value in {52/r=6,84/s=8,116/t=64,148/J=32}{
    \node[text=NNInput] at (\x,-12) {$\value$};
    \draw[NNLine,line width=.45pt] (\x-12,-20)--(\x+12,-20);
  }
  \node[ad derivative] at (52,-30) {$\bar r=8$};
  \node[ad derivative] at (84,-30) {$\bar s=8$};
  \node[ad derivative] at (116,-30) {$\bar t=\tfrac12$};
  \node[ad derivative] at (148,-30) {$\bar J=1$\\求导起点};
  \draw[ad reverse] (32,-47)--(22,-47);
  \node[anchor=west] at (35,-47) {梯度回传};
  \node[anchor=west,text=NNGradient] at (84,-47) {$\bar u=24+8=32$};
  \node[anchor=west] at (125,-47) {分支贡献相加};
\end{tikzpicture}

  \caption{同一算子图上的后向模式。节点位置与前向模式相同，前向数值保持不变，赭色虚线从$\bar J=1$开始回传敏感度；算子之间标出相应的局部导数。输入$u$收到经乘法传来的$24$和直接经加法传来的$8$，累加为$32$；输入$v$得到$16$。一次后向遍历同时得到两个偏导。}
  \label{fig:ffn-autodiff-reverse}
\end{figure*}

在图~\ref{fig:ffn-autodiff-reverse}中，$u$有两条通向目标的路径，所以$\bar u=24+8=32$。若漏掉直接进入加法的路径，就会错误地得到$24$。共享参数在多处被使用时也遵循同样规则：先保存各使用位置的局部回传贡献，再把它们累加到同一个参数。后向模式改变的是求导顺序，原函数的前向依赖并没有改变。

\FloatBarrier
\subsubsection{两种模式的计算代价}
\label{sec:ffn-autodiff-cost}

对一般映射$\bm{F}:\mathbb{R}^{p}\to\mathbb{R}^{q}$，记雅可比为$\bm{J}_{\bm{F}}\in\mathbb{R}^{q\times p}$。前向模式在给定输入方向$\bm{v}\in\mathbb{R}^{p}$时计算$\bm{J}_{\bm{F}}\bm{v}$，后向模式在给定输出权重$\bm{u}\in\mathbb{R}^{q}$时计算$\bm{J}_{\bm{F}}^{\top}\bm{u}$。二者都可直接计算乘积，无需存储完整雅可比。用逐坐标方向构造完整雅可比时，前向模式需要$p$个方向，后向模式需要$q$个方向；实际代价还取决于批处理、稀疏结构和基本算子的实现。

神经网络训练通常有很多参数，却只有一个标量目标，即$q=1$。后向模式从该目标出发，能够在与前向求值同阶的算术工作中得到完整参数梯度，代价是保存或重新计算反向所需的中间值。前述反向传播就是后向模式在神经网络损失上的具体应用\scite{baydin2018autodiff}

如果计算图包含多条损失分支，梯度在共同上游节点相加；如果某些分支被有意停止求导，它们就不再贡献上游梯度。离散采样、取整和按值选择分支也需要分别判断可微性：自动微分通常沿当前执行路径应用已登记的规则，不会自行补上对离散选择概率的导数。正确使用自动微分仍需理解目标、连接路径与不可微操作，梯度检查只是其中一种核验手段。

\section{理论分析}
\label{sec:ffn-theory}

一次预测可以计算、一个目标可以求导之后，还需回答两个不同的问题：这个网络族是否包含足够好的函数，以及有限样本和具体算法能否支持低预测风险。表达性考察前一个问题，可学习性考察后一个问题。把它们分开，才能解释为什么“有一组正确参数”不是“训练一定成功”的证明。

\subsection{表达性}
\label{sec:ffn-expressivity}

前章的异或例子是一个精确表示结论：指定输入集合上存在某组参数，输出规则与标签一致。面对一般连续输入，常见问题则是能否让网络与目标函数的误差足够小。\term{表达性}（expressivity）描述网络族能够表示或逼近哪些函数；它的判断对象是整类可选参数形成的函数，而不是某次训练得到的参数。

一个重要存在性结论是\term{通用逼近}（universal approximation）：允许适当增加网络规模时，可以逼近规定函数类中的任意目标。为避免把它理解成固定小网络的保证，下面明确给出目标域、误差尺度和允许变化的量。

\begin{theorem}[单隐藏层的通用逼近]
  \label{thm:ffn-universal-approximation}
  设$\bm{X}\subset\mathbb{R}^{d}$为非空紧集，$f:\bm{X}\to\mathbb{R}$连续，激活函数$\sigma:\mathbb{R}\to\mathbb{R}$连续且不是多项式。对于任意$\varepsilon>0$，存在有限宽度$m$及参数$a_j,b_j\in\mathbb{R}$、$\bm{w}_j\in\mathbb{R}^{d}$和$c\in\mathbb{R}$，使
  \begin{equation}
    \begin{aligned}
      g(\bm{x})&=c+\sum_{j=1}^{m}a_j\sigma(\bm{w}_j^{\top}\bm{x}+b_j),\\
      \sup_{\bm{x}\in\bm{X}}|g(\bm{x})-f(\bm{x})|&<\varepsilon\text{。}
    \end{aligned}
    \label{eq:ffn-universal-approximation}
  \end{equation}
\end{theorem}

这一定理是Leshno等关于非多项式激活结果在连续情形下的表述，隐藏层具有偏置，输出是隐藏响应的线性组合\scite{leshno1993approximation}Sigmoid、Tanh、ReLU、GELU与SiLU均满足这里的激活条件；ReLU不必处处可微，因为逼近定理并没有要求求导。向量值连续目标可以逐坐标应用该结论，再组合有限个隐藏单元。

定理中的量词顺序是先给定$f$与$\varepsilon$，再选择宽度和参数。它没有给出一个固定宽度对所有目标都够用的结论，也没有限制参数范数、保证梯度下降找到这些参数，或规定需要多少训练样本。误差控制只覆盖$\bm{X}$，不能据此保证域外外推。连续网络也不能在包含跳跃点的区间上，以任意小的一致误差逼近跳跃函数；改用按数据分布加权的误差，讨论对象与条件就会改变。

非线性本身也不是单隐藏层通用逼近的充分条件。例如，若激活为$\sigma(z)=z^2$，不论隐藏层多宽，线性输出都是输入坐标的次数至多为$2$的多项式。固定深度的多项式激活网络也只有有界次数，不能套用上述定理。这一反例针对规定深度和输出形式；若允许深度变化，函数类也会改变。

ReLU网络的一个直接读法是分段仿射变换。在一组隐藏单元的正负状态固定的输入区域内，所有ReLU都等价于保留某坐标或置零，整个网络因而是仿射的。不同激活状态形成不同区域，区域边界随参数改变。它能拼接出复杂形状，但固定有限网络仍只有有限个这样的区域。

一维情形能把这种拼接写成明确公式。设连续分段线性函数的内部转折点为$t_1<\cdots<t_m$，从左到右的斜率为$s_0,s_1,\ldots,s_m$，则可写成
\begin{equation}
  f(x)=c+s_0x+
  \sum_{j=1}^{m}(s_j-s_{j-1})\rho(x-t_j)\text{。}
  \label{eq:ffn-piecewise-linear-representation}
\end{equation}
每个ReLU在一个转折点之后改变斜率；常数$c$确定整体高度。若标准MLP不含直接的线性输入支路，可以用$x=\rho(x)-\rho(-x)$补出$s_0x$。因此宽度确实可以增加可表达的转折，而不只是抽象地增加参数。对紧区间上的连续函数取足够密的插值节点，再作分段线性插值，也给出了ReLU通用逼近在一维中的构造直觉。

深度提供另一种组织复杂函数的方式：重复复用已有变换。考虑$[0,1]$上的三角形函数
\begin{equation}
  T(x)=2\rho(x)-4\rho(x-1/2)+2\rho(x-1)\text{。}
  \label{eq:ffn-tent-function}
\end{equation}
它在$[0,1/2]$从$0$上升到$1$，在$[1/2,1]$再下降到$0$。令$T_k$表示$T$复合$k$次，而不是$T$的$k$次幂。$T_1$有两个线性区间，每个区间都把输入完整地送到$[0,1]$；再次应用$T$就把每个区间分成两段。归纳可得$T_k$在$[0,1]$有$2^k$个线性区间。

\begin{figure*}[htb]
  \centering
  % 原创解析示意。折点坐标由 T 的复合精确确定。
\begingroup
% 局部统一的科学绘图样式；曲线保持颜色和线型双重编码。
\pgfplotsset{
  ffn curve/.style={line width=1.3pt,line cap=round,line join=round},
  ffn reference curve/.style={NNInk!58,line width=0.95pt},
  ffn activation axis/.style={
    width=72mm,height=59mm,anchor=south west,
    axis lines=middle,
    axis line style={NNInk!65,line width=0.65pt},
    grid=major,grid style={NNLine!35,line width=0.3pt},
    tick align=outside,tick style={NNInk!55,line width=0.55pt},
    tick label style={font=\figurefont\footnotesize,text=NNInk},
    label style={font=\figurefont\footnotesize,text=NNInk},
    xlabel={$z$},ylabel={$\sigma(z)$},
    ylabel style={at={(axis cs:0,\pgfkeysvalueof{/pgfplots/ymax})},
      anchor=south,yshift=2pt},
    title style={font=\figurefont\footnotesize\bfseries,text=NNInk,
      at={(0,1.16)},anchor=south west},
    legend cell align=left,
    legend image post style={xscale=1.35},
    legend style={font=\figurefont\footnotesize,text=NNInk,draw=none,fill=none,
      at={(0.5,-0.18)},anchor=north,legend columns=1,
      inner sep=1pt,row sep=2pt,
      /tikz/every even column/.append style={column sep=3mm}},
    clip=true}}

\pgfplotsset{ffn depth axis/.style={
  width=50mm,height=43mm,anchor=south west,
  xmin=0,xmax=1,ymin=0,ymax=1.1,xtick={0,0.5,1},ytick={0,1},
  axis lines=left,axis line style={NNInk!65,line width=0.65pt},
  grid=major,grid style={NNLine!35,line width=0.3pt},
  tick align=outside,tick style={NNInk!55,line width=0.55pt},
  tick label style={font=\figurefont\footnotesize,text=NNInk},
  label style={font=\figurefont\footnotesize,text=NNInk},xlabel={$x$},
  title style={font=\figurefont\small,text=NNInk},clip=false}}
\begin{tikzpicture}[font=\figurefont\footnotesize]
  \begin{axis}[ffn depth axis,at={(0,0)},title={$T_1(x)$}]
    \addplot[NNInput,ffn curve] coordinates {(0,0) (0.5,1) (1,0)};
    \node[font=\figurefont\footnotesize,text=NNInput,anchor=north]
      at (axis description cs:0.5,-0.32) {2个线性区间};
  \end{axis}
  \begin{axis}[ffn depth axis,at={(56mm,0)},title={$T_2(x)$}]
    \addplot[NNHidden,ffn curve]
      coordinates {(0,0) (0.25,1) (0.5,0) (0.75,1) (1,0)};
    \node[font=\figurefont\footnotesize,text=NNHidden,anchor=north]
      at (axis description cs:0.5,-0.32) {4个线性区间};
  \end{axis}
  \begin{axis}[ffn depth axis,at={(112mm,0)},title={$T_3(x)$}]
    \addplot[NNOutput,ffn curve]
      coordinates {(0,0) (0.125,1) (0.25,0) (0.375,1) (0.5,0)
        (0.625,1) (0.75,0) (0.875,1) (1,0)};
    \node[font=\figurefont\footnotesize,text=NNOutput,anchor=north]
      at (axis description cs:0.5,-0.32) {8个线性区间};
  \end{axis}
\end{tikzpicture}
\endgroup

  \caption{三角形函数的复合产生更多转折。三幅解析图分别为$T_1=T$、$T_2=T\circ T$和$T_3=T\circ T\circ T$，横轴均为$x\in[0,1]$，纵轴为函数值。每次复合将线性区间数翻倍；这是一维构造，不是实际训练实验，也不表示任意目标都能从增加深度中获得相同收益。}
  \label{fig:ffn-depth-composition}
\end{figure*}

把$T$的三个ReLU响应与线性组合逐次串联，可以用深度随$k$线性增长、每层宽度有常数上界的网络表示$T_k$；相邻线性映射可合并到下一层权重中。单隐藏层的一维ReLU网络中，每个单元至多引入一个内部转折，所以要精确表示具有$2^k$个线性区间的$T_k$，至少需要$2^k-1$个隐藏单元。图~\ref{fig:ffn-depth-composition}展示了这种重复组合的效率。更一般的深度分离结果还需固定误差度量、目标函数族和比较网络的范围\scite{telgarsky2016depth}

这个例子比较的是同一目标的表示成本，并没有证明深度越大就越容易训练或越能泛化。某些目标用浅层结构已经足够；若目标本身具有可复用的组合规律，深度才可能以较少单元组织这种规律。表达能力的分析因此应同时说明目标、激活、深度、宽度和逼近精度。

\subsection{可学习性}
\label{sec:ffn-learnability}

通用逼近允许观察目标函数之后选择参数，而训练只能看到有限个样本。网络可以在这些样本之间构造许多不同的延伸，未见输入上的表现也就可能不同。\term{可学习性}考察给定任务、数据假设和学习方式时，样本能否支持低风险的预测规则。统计上需要多少样本、计算上能否找到合适参数，分别需要复杂度和优化方面的依据。

沿用统计学习框架，假设训练样本独立同分布地来自$\mathcal{P}$，测试样本也来自同一分布。对事先确定的网络函数空间$\mathcal{H}$，记期望风险为$R(f)=\mathbb{E}_{\mathcal{P}}\ell(f(\bm{x}),y)$、经验风险为$\widehat R_S(f)$，学习算法的输出为$\widehat f=A(S)\in\mathcal{H}$。若对所有$f\in\mathcal{H}$同时有$|R(f)-\widehat R_S(f)|\leq\Delta$，且算法的经验优化误差至多为$\xi$，则
\begin{equation}
  \begin{aligned}
    R(\widehat f)-R^*\leq{}&
    \underbrace{\inf_{f\in\mathcal{H}}R(f)-R^*}_{\text{逼近误差}}\\
    &+\underbrace{2\Delta}_{\text{估计误差的控制}}
    +\underbrace{\xi}_{\text{优化误差}}\text{。}
  \end{aligned}
  \label{eq:ffn-risk-decomposition}
\end{equation}
这里$R^*$是同一损失下允许所有可测预测规则得到的最优风险，$\xi$满足$\widehat R_S(\widehat f)\leq\inf_{f\in\mathcal{H}}\widehat R_S(f)+\xi$。把$\widehat f$的经验风险与某个近似最优比较函数的经验风险连接起来，再分别用一致偏差$\Delta$替换为期望风险，就得到该式。对于无界损失，要建立这样的统一界还需矩或尾部等额外条件。

这三项解释了增加模型规模的不同作用：更大的函数空间可能降低逼近误差，却会改变估计误差的控制；更深的计算也可能使优化误差难以降低。式\eqref{eq:ffn-risk-decomposition}是一个有条件的风险分解，不保证某个具体界在大网络中紧，也不意味着真实测试误差必须随参数量单调增加。

输入维度通过表示规模与数据几何共同起作用。若每个坐标都满足$|x_j|\leq c$，只能粗略保证$\lVert\bm{x}\rVert_2\leq c\sqrt d$，因而以输入范数为依据的界可能随$d$变大；若输入整体已归一化，或数据只占据较低维的有效区域，维度的作用就不能只按坐标数判断。但低维分布本身也不保证任务简单：标签可能在该区域内变化得很复杂，还可能含有噪声。输入分布、目标结构和损失应在同一分析中固定。

样本量的作用可以用前文已经证明的两层ReLU结论具体展示。此处改用标签$y\in\{-1,+1\}$和标量分类得分，讨论不含偏置的函数
\[
  f(\bm{x})=\sum_{j=1}^{m}a_j\rho(\bm{w}_j^{\top}\bm{x}),
  \qquad \sum_{j=1}^{m}|a_j|\leq A,
  \qquad \lVert\bm{w}_j\rVert_2\leq B\text{。}
\]
在抽样前固定$A,B,r,\gamma>0$，并假定$\lVert\bm{x}\rVert_2\leq r$。由
\BookTheoremRef{thm:two-layer-norm-risk}{第一卷“基础、理论和可信性”的“深度学习的可学习性”章“两层网络的范数与期望风险”定理}，
以至少$1-\delta$的概率，上述函数同时满足
\begin{equation}
  \begin{aligned}
    \mathbb{P}_{\mathcal{P}}\{yf(\bm{x})\leq0\}
    \leq{}&\frac1n\sum_{i=1}^{n}
      \mathbb{I}\{y_i f(\bm{x}_i)\leq\gamma\}\\
    &+\frac{4ABr}{\gamma\sqrt n}
      +3\sqrt{\frac{\log(2/\delta)}{2n}}\text{。}
  \end{aligned}
  \label{eq:ffn-margin-risk}
\end{equation}
其中$\gamma$是要求训练样本达到的正分类间隔，$\delta\in(0,1)$是允许保证失效的概率。左侧把零间隔也计为错误，因而可控制固定平局规则下的分类风险。右侧既包含未达到间隔的训练样本比例，也包含从有限样本推断新样本时的复杂度代价；它不只检查训练样本是否分对。

在$A,B,r,\gamma,\delta$固定时，将样本量增至四倍，会把后两项减半。这是该充分界的缩放关系，不意味着实际测试错误率一定减半。它也说明为什么要区分独立样本与重复记录：把同一条记录复制四次没有产生这项独立同分布保证所要求的新信息。

参数范数与间隔需要一起看。把所有输出权重乘以一个很小的正数，会缩小$A$，但也按相同比例缩小分类间隔，分类符号并不改变；因此不能由权重变小直接推出分类风险变小。若在保持函数不变的前提下，将某个ReLU隐藏单元的输入权重乘$c>0$、输出权重除以$c$，不同层的原始范数同样会变化。分析应采用与函数、间隔及参数化相配的复杂度量。

式\eqref{eq:ffn-margin-risk}没有显式出现隐藏宽度$m$，因为总输出权重和各隐藏权重范数已经受到限制。增加宽度可以提供更多候选特征，却不能无限放大所有特征。如果宽度增加时$A$、$B$也增加，或有效间隔变小，界仍会变差。允许所有宽度和任意权重的分类函数空间又可能具有无限VC维，不能把这条有范数与间隔条件的保证推广为整个网络族的分布无关PAC保证。固定有限ReLU架构的VC分析及其统计学习结论见
\BookChapterRef{chap:deep-learning-theory}{第一卷“基础、理论和可信性”的“深度学习的可学习性”章}。

深度影响层间扰动的放大。对本章的串联网络，若隐藏激活均为$1$-Lipschitz，即不会增大两点之间的欧氏距离，原始输出映射满足
\begin{equation}
  \lVert\bm{a}(\bm{x})-\bm{a}(\bm{x}')\rVert_2
  \leq
  \left(\prod_{\ell=1}^{L+1}\lVert\bm{W}_\ell\rVert_2\right)
  \lVert\bm{x}-\bm{x}'\rVert_2\text{。}
  \label{eq:ffn-input-sensitivity}
\end{equation}
这里矩阵谱范数$\lVert\bm{W}_\ell\rVert_2$表示线性变换对向量长度的最大放大倍数。每层分别使用这一上界，再沿函数复合相乘，便得到该式；有其他Lipschitz常数的激活还需乘上相应因子。偏置在比较同一层的两个输入时相消，但它仍会影响激活区域及函数空间的复杂度。

式\eqref{eq:ffn-input-sensitivity}控制固定网络的输入敏感度，不能单独给出泛化保证。深层谱复杂度分析还结合各层其他范数、输入尺度和分类间隔，具体条件见
\BookSectionRef{sec:network-norm-bounds}{第一卷“基础、理论和可信性”的“深度学习的可学习性”章“基于参数大小的可学习性”节}。
它说明深度不能只作为层数进入判断：层间尺度的乘积、可复用目标结构以及训练能否得到合适表示，都可能同时变化。

学习算法决定从许多低训练损失参数中选择哪一个解。
\BookSectionRef{sec:implicit-bias}{第一卷“基础、理论和可信性”的“深度学习的可学习性”章“基于参数梯度更新的可学习性”节}
讨论的选解偏好，
\BookSectionRef{sec:neural-tangent-kernel}{第一卷“基础、理论和可信性”的“深度学习的可学习性”章“基于参数梯度初始值的可学习性”节}
讨论的初始梯度特征，分别在规定条件下描述这种选择。初始化、更新方式、正则化和停止时刻既可能影响经验优化误差，也可能影响输出对样本替换的敏感性。若算法具有损失均匀稳定性$\beta_n$，其期望泛化差可受$\beta_n$控制；这条结论比较的是更换训练样本后的两个训练结果，不能用同一次训练中两次参数更新很接近来代替。

因此，从前馈结构得到学习判断，需要把任务、样本分布、候选函数约束和算法放在一起。独立同分布风险界不直接覆盖部署分布改变，交叉熵降低也不直接等于零一风险保证。一个完整论证应当说明网络能逼近什么、样本支持多强的估计、算法能将经验目标降到哪里，以及这些条件如何共同转成所关心的预测风险。

\section{常见应用}
\label{sec:ffn-applications}

前馈网络既可以直接完成固定长度输入的预测，也可以在更大的模型中承担一段变换。应用时需要明确它接收什么表示、输出什么对象、由哪个损失训练。仅给出隐藏层的形状，还不能得到完整的预测方法。

\subsection{表格数据的分类与回归}
\label{sec:ffn-tabular}

表格中的一行通常描述一个对象，不同列分别记录数值或类别属性。网络接收的是定长数值向量，所以必须先规定每列怎样进入向量。数值列可按训练集统计量标准化，类别列可编码成互斥的指示坐标，缺失值则需规定填充值并酌情加入缺失标记。这些规则也属于模型使用过程：验证集、测试集和新样本使用训练阶段已经确定的转换，不重新用测试数据估计均值和方差。

考虑一个教学任务：由设备使用时长、近期告警次数和三个设备类别之一，预测下一段时间是否发生故障。设两个数值变量标准化为$\widetilde x_1,\widetilde x_2$，设备类别编码为$\bm{c}\in\{0,1\}^{3}$且坐标和为$1$，输入为
\[
  \bm{x}=(\widetilde x_1,\widetilde x_2,c_1,c_2,c_3)^{\top}
  \in\mathbb{R}^{5}\text{。}
\]
采用宽度为$16$和$8$的两个ReLU隐藏层、一个标量输出，可以把完整模型写为
\begin{equation}
  \begin{aligned}
    \bm{h}^{(1)}&=\rho(\bm{W}_1\bm{x}+\bm{b}_1),\\
    \bm{h}^{(2)}&=\rho(\bm{W}_2\bm{h}^{(1)}+\bm{b}_2),\\
    p&=s(\bm{w}_3^{\top}\bm{h}^{(2)}+b_3)\text{。}
  \end{aligned}
  \label{eq:ffn-tabular-classifier}
\end{equation}
其中$\bm{W}_1\in\mathbb{R}^{16\times5}$、$\bm{W}_2\in\mathbb{R}^{8\times16}$、$\bm{w}_3\in\mathbb{R}^{8}$，偏置维度与对应输出相同。这组层宽是示例设计，不是经过实验验证的推荐配置。

训练标签$y\in\{0,1\}$记录指定时间窗内是否发生故障，用二元交叉熵构造式\eqref{eq:ffn-training-objective}，再按算法~\ref{alg:ffn-minibatch-learning}学习参数。输入必须只包含预测时刻已知的信息；故障发生后记录的维修结果不能作为预测前的输入。若同一设备有多条记录，还需按设备和时间等关系划分数据，使评价符合预期部署场景。

预测新设备记录时，先按保存的规则得到$\bm{x}$，再计算$p$，根据验证阶段确定的决策阈值得到类别。若任务改为预测后续维修时长$y\in\mathbb{R}$，可保留输入编码与隐藏层，改用线性输出和平方损失；若必须保证非负预测，可增加相应输出约束。改变输出和损失是在改变预测对象，不能只把故障标签换成一个数值而保留原来的概率解释。

全连接结构允许数值与类别特征在隐藏层共同作用，但它没有自动规定某种单调关系，也没有保证类别编码在缺少样本时得到可靠估计。类别很多时，直接增加指示坐标还会扩大第一层参数量，可以考虑用可学习的低维向量表示类别，但这种表示同样需要训练数据支持。对表格任务，应与适合任务的简单模型比较，并依据独立评价选择模型；架构名称不能替代数据与误差分析。

\subsection{其他网络中的前馈模块}
\label{sec:ffn-modules}

当输入表示已经由其他模型得到时，前馈网络仍可以把它投影到新空间，或转换成任务输出。\term{特征投影}（feature projection）改变表示的维度或坐标组织，例如用$\bm{W}\bm{h}+\bm{b}$把$d$维向量映射到$q$维。若只是两次线性投影且中间没有非线性，它们仍可以合并为一次仿射映射；如果中间维度较小，这种参数化还会限制整体线性映射的秩。

\term{分类头}（classification head）是把已有表示转换为类别得分的末端模块。卷积编码器可以把图像压成一个向量，再接线性层或MLP得到类别得分。联合训练时，分类损失的梯度沿前馈模块继续传到编码器；若固定编码器，只更新分类头，学习的参数范围就缩小了。同一个模块形状可以对应不同训练过程，应明确哪些参数参与更新。

有时模块输出的是两个对象的相关性或匹配程度。例如，将两个$d$维表示$\bm{u},\bm{v}$拼接，构成输入$[\bm{u};\bm{v}]\in\mathbb{R}^{2d}$，再计算
\begin{equation}
  \operatorname{score}(\bm{u},\bm{v})=
  \bm{w}^{\top}\sigma\!\left(\bm{W}[\bm{u};\bm{v}]+\bm{b}\right)+c\text{。}
  \label{eq:ffn-scoring-module}
\end{equation}
其中隐藏宽度为$m$，$\bm{W}\in\mathbb{R}^{m\times2d}$、$\bm{b},\bm{w}\in\mathbb{R}^{m}$、$c\in\mathbb{R}$。这个\term{评分函数}（scoring function）把对象表示转换为标量得分，概率变换和损失由匹配或排序任务决定。拼接顺序意味着交换两个输入可能改变得分；若任务要求对称性，需要在输入或计算结构中加入相应约束。

前馈模块也可以对一组位置重复使用。设第$t$个位置的表示为$\bm{h}_t\in\mathbb{R}^{d}$，用同一组参数计算
\begin{equation}
  \operatorname{FFN}(\bm{h}_t)=
  \bm{W}_2\sigma(\bm{W}_1\bm{h}_t+\bm{b}_1)+\bm{b}_2\text{，}
  \label{eq:ffn-positionwise-module}
\end{equation}
其中$\bm{W}_1\in\mathbb{R}^{m\times d}$、$\bm{W}_2\in\mathbb{R}^{d\times m}$，偏置维度与相应输出一致。这种\term{逐位置变换}（position-wise transformation）分别改变每个位置内部的特征组合，在Transformer等模型中作为一个模块使用。它读取的$\bm{h}_t$可以已经包含其他位置的信息，但这次变换本身不读取$\bm{h}_{t'}$，因而不直接建立新的跨位置交互。

位置共享参数意味着所有位置的梯度都要累积到同一组$\bm{W}_1,\bm{W}_2$，而不是为每个位置保存一份独立参数。这个例子也区分了网络名称与模块职责：更大的模型可以包含前馈计算，前馈模块负责特征变换，其他结构负责空间、序列或对象之间的信息交换。后续各模型章会在这一基础上讨论相应的结构假设。

\section{本章小结}
\label{sec:ffn-summary}

前馈神经网络把可学习表示落实为一条没有有向环的计算路径。全连接层组合已有坐标，非线性激活让多层复合超出仿射映射，输出变换与损失共同规定训练所要预测的对象。选择层宽和激活时，既要看能够形成怎样的函数，也要看参数量、计算量和梯度传播条件。

隐藏层的学习信号来自最终目标对它的敏感度。反向传播按链式法则逐层汇总这种影响，用局部导数和已有活动值计算所有参数梯度；自动微分将同样的原则推广到程序的实际计算路径。初始化决定起点，优化器使用梯度作出更新，二者都影响学习结果，正确求导只是完成训练的一个条件。

能够逼近目标函数、能够优化经验目标、能够用有限样本控制预测风险，仍是三个不同的要求。通用逼近解释存在怎样的参数，深度与宽度解释表示成本，复杂度和稳定性则在数据与算法条件下连接训练结果和新样本风险。前馈网络适合处理定长表示，也可作为投影、分类、评分和逐位置变换模块；遇到具有明确空间、序列或关系结构的数据时，还需判断全连接假设是否充分利用了这些规律。后续章节将在相应结构上继续建立模型。
